\providecommand{\srcrem}{\mathrm{rem}}
\providecommand{\srcB}{\mathcal{B}}
\providecommand{\srcT}{\mathcal{T}}
\providecommand{\srcRR}{\mathcal{R}}
\providecommand{\srctr}{\mathsf{T}}
\section{Spanning, completion and integral frames}\label{sec:spn}

\providecommand{\spnHS}{\mathrm{HS}}
\providecommand{\spncodim}{\operatorname{codim}}
\providecommand{\spnim}{\operatorname{im}}
\providecommand{\spnE}{\mathbb{E}}
\providecommand{\spnP}{\mathbb{P}}
\providecommand{\spnTV}{d_{\mathrm{TV}}}
\providecommand{\spnStab}{\operatorname{Stab}}
\providecommand{\spnMat}{\operatorname{Mat}}
\providecommand{\spnDD}{\mathfrak{D}}
\providecommand{\spnT}{\mathsf{T}}
\providecommand{\spnqbin}[3]{\genfrac{[}{]}{0pt}{}{#1}{#2}_{#3}}
\providecommand{\spnOrth}{\mathcal{O}}
\providecommand{\spnTame}{\mathcal{T}}
\providecommand{\spnEmb}{\mathcal{E}}
\providecommand{\spnAlt}{\operatorname{Alt}}
\providecommand{\spncoker}{\operatorname{coker}}

The source proof (\cref{src:sec:mechanism}) compares an upper and a lower bound
for one twisted lattice, and its determinant comparison
(\cref{src:sec:comparison}) needs a short family that spans almost all of the
target, together with integral vectors completing it to a basis. This section
treats that spanning and completion step and its integral refinements: the
second-moment spanning lemma at sharp constants and at the source's own
constants, exact finite selection certificates, why near-full spanning is
needed and how missing dimensions can be paid for instead, optimal integral
tight frames and self-dual projective lattices, and lower bounds for integral
Gram determinants with vanishing minors. Every result is \stU. Notation is local
to each subsection. The random matrix compressions controlled by the
positive-type lemma of the source's spanning step are in \cref{sec:cmp}, the
auxiliary anticyclotomic characters and their conductors in \cref{sec:cost}
(\cref{spn:sec:characters,spn:sec:conductor,spn:sec:benchmark}), and the
annihilator of the odd cyclotomic cocycles in \cref{sec:sel}
(\cref{spn:sec:annihilator}).

%======================================================================
\subsection{A sharp spanning interface}\label{spn:sec:spanning}
\sources{grh_spanning_and_s5_2026-11-26; grh_path_rigidity_2026-11-27}

Let $S,T$ be finite-dimensional complex Hilbert spaces, with $s=\dim S>0$. Let a
finite group $G$ act unitarily on $T$, by $\rho$, and let $P$ be the orthogonal
projection onto a $G$-stable subspace $W$, of dimension $d>0$. Thus
$P\rho(g)=\rho(g)P$. Fix a linear map $B\colon S\to T$ and write
$C_g=P\rho(g)B\colon S\to W$.

The spanning lemma of the source, \code{isSpread_of_second_moment}, rests on
two quantitative hypotheses: for every subspace $U\le S$,
\begin{equation}\label{spn:eq:near}
 \norm{PBP_U}_{\spnHS}^2\ge \dim U-us,
\end{equation}
and
\begin{equation}\label{spn:eq:moment}
 \frac d{\abs{G}}\sum_{g\in G}\norm{B^*\rho(g)B}_{\spnHS}^2\le(1+v)s^2 .
\end{equation}
Here $P_U$ denotes orthogonal projection onto $U$, and $0\le u<1$, $v\ge0$.
Condition \eqref{spn:eq:near} is the predicate \code{IsNearIsometry} of the
source; it is a statement on every source subspace, not merely one lower
bound on the total Hilbert--Schmidt norm.

\begin{lemma}[The square-root estimate; \stU]\label{spn:lem:sqrt}
For every subspace $F\le W$, of relative dimension $\alpha=\dim F/d$, some
$g\in G$ satisfies
\begin{equation}\label{spn:eq:sqrt}
 \dim C_g^{-1}(F)\le s\bigl[u+\sqrt{(1+v)\alpha}\bigr].
\end{equation}
In particular, $s\le d/(1-u)$.
\end{lemma}

\begin{proof}
Put $M=\abs{G}^{-1}\sum_g\rho(g)BB^*\rho(g)^*$. Unitarity and a change of
variables in the double sum give
\[
 \tr(M^2)=\frac1{\abs{G}}\sum_g\norm{B^*\rho(g)B}_{\spnHS}^2 .
\]
The Hilbert--Schmidt Cauchy--Schwarz inequality therefore gives
\[
 \frac1{\abs{G}}\sum_g\norm{P_F\rho(g)B}_{\spnHS}^2=\tr(P_FM)
 \le\sqrt{\dim F\,\tr(M^2)}\le s\sqrt{\alpha(1+v)} .
\]
Choose $g$ with its summand at most this average and put $U=C_g^{-1}(F)$. On
$U$, the maps $C_g$ and $P_F\rho(g)B$ agree. Because $C_g=\rho(g)PB$, and
right-composition by an orthogonal projection does not increase the
Hilbert--Schmidt norm, \eqref{spn:eq:near} yields
\[
 \dim U-us\le\norm{PBP_U}_{\spnHS}^2=\norm{C_gP_U}_{\spnHS}^2
 \le\norm{P_F\rho(g)B}_{\spnHS}^2 .
\]
This proves \eqref{spn:eq:sqrt}. For the last assertion, apply
\eqref{spn:eq:near} to $\ker(PB)$. It gives $\dim\ker(PB)\le us$, hence
$(1-u)s\le\rank(PB)\le d$.
\end{proof}

The source first replaces both errors by a common $\eps$, then applies
$\sqrt{\alpha(1+\eps)}\le(\alpha+1+\eps)/2$ before proceeding to the greedy
argument. Keeping the square root retains the information near the limiting
dimension. Set
\begin{equation}\label{spn:eq:parameters}
\begin{gathered}
 a=1-u,\qquad z=1+v,\qquad \alpha_*=\frac{a^2}{z},\\[3pt]
 \delta_*=1-\alpha_*=\frac{2u+v-u^2}{1+v},\qquad h=\frac{sz}{2ad},\qquad
 q=\max\set{0,1-h}.
\end{gathered}
\end{equation}
Thus $a,z,\alpha_*,h>0$, $0\le\delta_*<1$, and $0\le q<1$.

\begin{theorem}[Greedy selection with the sharp limiting loss; \stU]\label{spn:thm:greedy}
Assume \eqref{spn:eq:near} and \eqref{spn:eq:moment}. For every integer $t\ge0$
there is a set $A\subseteq G$, with $\abs{A}\le t$, such that
\begin{equation}\label{spn:eq:greedy}
 \spncodim_W\sum_{g\in A}\spnim C_g\le d\bigl[\delta_*+\alpha_*q^t\bigr].
\end{equation}
Consequently, for $0<\tau<\alpha_*$, a loss of at most $(\delta_*+\tau)d$ is
achieved with at most
\begin{equation}\label{spn:eq:count}
 \ceil*{\frac{2ad}{sz}\log\frac{\alpha_*}{\tau}}
\end{equation}
translates. The span of all the images has codimension at most
$\floor{\delta_*d}$.
\end{theorem}

\begin{proof}
Let the current span be $F$, with codimension $c$. Rank--nullity for
$S\to W/F$ gives
\[
 \dim(F+\spnim C_g)-\dim F=s-\dim C_g^{-1}(F).
\]
By \cref{spn:lem:sqrt}, some $g$ therefore adds at least
\begin{equation}\label{spn:eq:gain}
 s\Bigl[a-\sqrt{z(1-c/d)}\Bigr]
\end{equation}
dimensions. Suppose $c>\delta_*d$, and put $e=c-\delta_*d>0$. The bracket in
\eqref{spn:eq:gain} is
\[
 a-\sqrt{a^2-ze/d}=\frac{ze/d}{a+\sqrt{a^2-ze/d}}\ge\frac{ze}{2ad}.
\]
The square root is real because $c\le d$. Thus, after the selected step, the
new excess codimension is at most $(1-h)e$, and hence at most $qe$. If
$c\le\delta_*d$, keep the existing set of indices. Induction from the empty
span, $c_0=d=\delta_*d+\alpha_*d$, proves \eqref{spn:eq:greedy}. Taking a set of
indices accommodates repetitions automatically.

Since $q\le e^{-h}$, the number in \eqref{spn:eq:count} suffices. This also
covers $q=0$, as the displayed ceiling is at least one. Finally, take
$F=\sum_{g\in G}\spnim C_g$ in \eqref{spn:eq:sqrt}. Every preimage is all of
$S$, so $1\le u+\sqrt{z\dim F/d}$, whence $\dim F\ge\alpha_*d$. Integrality
gives the stated floor.
\end{proof}

With $d$ and $s$ fixed, the possible total spans are determined exactly. Put
$F=\sum_{g\in G}\spnim C_g$ and $r=\dim F$; the spaces, the group, the representation and the map
may vary.

\begin{theorem}[The complete feasible rank spectrum; \stU]\label{spn:thm:rank-spectrum}
For fixed $(d,s,u,v)$ with $0\le u<1$ and $v\ge0$, a system satisfying \eqref{spn:eq:near} and
\eqref{spn:eq:moment} exists if and only if $as\le d$, and then the possible values of $r$ are exactly
\[
 \set{r_*,r_*+1,\dots,d},\qquad r_*=\max\Bigl\{\ceil{as},\ceil*{\frac{a^2d}z}\Bigr\}.
\]
Every such $r$ is attained with $T=W=\C^d$, $P=I$, a cyclic group of order $r$, a contraction $B$
and $\rank B=\ceil{as}$.
\end{theorem}

\begin{proof}
\Cref{spn:lem:sqrt} gives $r\ge\rank(PB)\ge as$, so $as\le d$ is necessary. With $M$ as in its proof
let $A=(PMP)|_W$; since $P$ commutes with $\rho$, $A$ is the covariance of the maps $C_g$, its image
is $F$, and $\tr A=\norm{PB}_{\spnHS}^2\ge as$. In the decomposition $T=W\oplus W^\perp$,
$\tr(M^2)\ge\tr(A^2)$, so Cauchy--Schwarz over the $r$ positive eigenvalues of $A$ gives
$(as)^2\le r\tr(A^2)\le r\tr(M^2)\le rzs^2/d$, that is $r\ge a^2d/z$.

For the construction let $r_*\le r\le d$ and $k=\ceil{as}$, so $1\le k\le\min(s,r)$. Choose
$\lambda_1,\dots,\lambda_k\in(0,1]$ with sum $as$ (the first $k-1$ equal to one), let
$E\colon\C^s\to\C^r$ send $e_j$ to $\sqrt{\lambda_j}e_j$ for $j\le k$ and to $0$ otherwise, let $Q$ be
the unitary Fourier matrix of order $r$, and put $B=QE$, embedded in the first $r$ coordinates of
$T=W=\C^d$, with $P=I$. As $B^*B\le I_S$ and $\tr(I_S-B^*B)=us$, positivity gives
$\norm{BP_U}_{\spnHS}^2=\dim U-\tr(P_U(I_S-B^*B))\ge\dim U-us$ for every $U$. Let the generator of
$\Z/r\Z$ act by $\operatorname{diag}(1,\omega,\dots,\omega^{r-1})$, $\omega=e^{2\pi i/r}$, on the first $r$
coordinates and trivially on the others. Distinct Fourier characters kill the off-diagonal entries
of the averaged covariance, and every entry of $Q$ has modulus $r^{-1/2}$, so $M=(as/r)P_{\C^r}$: the
orbit images span the first $r$ coordinates, and
$\frac d{\abs G}\sum_g\norm{B^*\rho(g)B}_{\spnHS}^2=d\tr(M^2)=d(as)^2/r\le zs^2$. For $r=1$ use the
trivial group. Finally $a^2d/z\le d$.
\end{proof}

\begin{proposition}[Equality in the second-moment bound; \stU]\label{spn:prop:rank-equality}
A feasible system has $r=a^2d/z$ if and only if $\spnim B\subseteq W$, $0\le B^*B\le I_S$ with
$\tr(B^*B)=as$, and $M=(as/r)P_F$. Such a system exists exactly when $a^2d/z$ is an integer at least
$\ceil{as}$.
\end{proposition}

\begin{proof}
Equality throughout the chain above forces the off-diagonal and $W^\perp$ blocks of $M$ to vanish, so
$0=\abs G^{-1}\sum_g\norm{(I-P)\rho(g)B}_{\spnHS}^2=\norm{(I-P)B}_{\spnHS}^2$; it makes the positive
eigenvalues of $A$ equal and their sum $as$, whence $M=(as/r)P_F$ and $\tr(B^*B)=as$. Applying
\eqref{spn:eq:near} to the span of the eigenvectors of $B^*B$ with eigenvalue $\mu_j<1$ gives
$\sum_j(1-\mu_j)_+\le us=\sum_j(1-\mu_j)$, so every $\mu_j\le1$. Conversely, a contraction of trace
$as$ satisfies \eqref{spn:eq:near} by the positivity argument, and the covariance gives equality in
\eqref{spn:eq:moment}. Existence follows from \cref{spn:thm:rank-spectrum}.
\end{proof}

There are two independent obstructions: one map needs $\ceil{(1-u)s}$ image directions, and the
second moment forces $\ceil{(1-u)^2d/(1+v)}$ orbit directions; the larger is attained, and so is every
larger rank. For $(d,s,u,v)=(100,90,\tfrac1{10},\tfrac15)$ the two bounds are $81$ and $68$ and the
answer is $81$; for $(100,10,\tfrac1{10},\tfrac15)$ they are $9$ and $68$ and the answer is $68$; for
$(100,120,\tfrac1{10},\tfrac15)$ there is no system, as $(1-u)s=108>d$. Exact rational checks of the
thresholds and constructions are in
\artifact{grh_path_rigidity_2026-11-27/checks/rank_checks.json}.

\begin{corollary}[Sharpness under the moment hypotheses; \stU]\label{spn:prop:sharp}
For every rational pair $0\le u<1$, $v\ge0$, there are examples satisfying
\eqref{spn:eq:near} and \eqref{spn:eq:moment} whose total span has dimension
exactly $\alpha_*d$. Thus $\delta_*$ cannot be decreased as a universal bound
under these hypotheses.
\end{corollary}

\begin{proof}
Choose $d$ with $\alpha_*d$ a positive integer and take $s=1$; then $r_*=\alpha_*d$ in
\cref{spn:thm:rank-spectrum}. (Equivalently, $B(1)=\sqrt{1-u}\,e_1$ with the cyclic group of order
$r=\alpha_*d$ rotating the first $r$ coordinates.)
\end{proof}

\begin{corollary}[Equal errors]\label{spn:cor:equal}
When $u=v=\eps$, the sharp limiting deficit is
$\delta_*=(3\eps-\eps^2)/(1+\eps)$. For $\eps=p/v_0$, with integers
$0\le p<v_0$, one explicit equality example has $d=v_0(v_0+p)$ and
$r=(v_0-p)^2$.
\end{corollary}

The total-rank argument is an instance of the classical second-moment, or
effective-rank, inequality \cite{Welch1974Lower,Mixon2020Three}. Greedy rank
accumulation is likewise a classical technique; see the submodular-cover
analysis in \cite{Wolsey1982Analysis}. What is extracted here is the exact
threshold $\delta_*$ of the source hypotheses and an explicit selection
schedule.

\begin{lemma}[Restriction loss; \stU]\label{spn:lem:restriction}
Let $S_0\le S$ have codimension $r_0$. The same selected set in
\cref{spn:thm:greedy} satisfies
\begin{equation}\label{spn:eq:restricted}
 \spncodim_W\sum_{g\in A}C_g(S_0)\le d\bigl[\delta_*+\alpha_*q^t\bigr]+tr_0 .
\end{equation}
\end{lemma}

\begin{proof}
The direct sum of $\abs{A}$ copies of $S/S_0$ surjects onto
$\bigl(\sum_{g\in A}C_g(S)\bigr)/\bigl(\sum_{g\in A}C_g(S_0)\bigr)$. Thus the
difference of the two span dimensions is at most $\abs{A}r_0\le tr_0$. Apply
\eqref{spn:eq:greedy}.
\end{proof}

Inequality \eqref{spn:eq:restricted} balances the benefit of selecting more
translates against the cost of restricting the source. If $r_0/s=\theta>0$ and
$2\theta<a$, taking $t=\ceil{h^{-1}\log(a/(2\theta))}$ gives
\begin{equation}\label{spn:eq:optimized}
 \frac1d\spncodim_W\sum_{g\in A}C_g(S_0)
 \le\delta_*+\frac{2a\theta}{z}\Bigl(1+\log\frac{a}{2\theta}\Bigr)
 +\frac{\theta s}{d}
 \le\delta_*+\frac{2a\theta}{z}\Bigl(1+\log\frac{a}{2\theta}\Bigr)
 +\frac{\theta}{a}.
\end{equation}
Indeed, $d\alpha_*h=as/2$, and rounding up increases the linear loss by at most
$r_0$.

%======================================================================
\subsection{The source spanning theorem at its original constants}\label{spn:sec:source-improvement}
\sources{grh_spanning_and_s5_2026-11-26}

The parameter called $b$ here is the large weight in the source, not the moment
factor $z=1+v$. The raw estimates in the source file
\code{SpanningTheorem.lean} are
\[
 u=\frac5{b^4},\qquad v=\frac9{b^4};
\]
they occur in the following declarations, respectively:
\begin{itemize}
\item \code{isNearIsometry_e}\textsubscript{0}\code{_comp_scaledPoincare};
\item \code{targetDim_mul_sum_hilbertSchmidtNorm_sq_le}.
\end{itemize}
The proof of \code{isSpread_goodPoincare} later replaces both by
$\eps=(8b^2)^{-1}$; that enlargement is unnecessary for the argument here.
The source-restriction argument in \code{ShortAndEnough.lean} has
$r_0/s\le b^{-8}$. The maps are identified through \code{sourceFormCoords};
the nonzero scalar normalizations in \code{goodPoincare} and
\code{normalizedPoincare} do not change the image subspaces.

\begin{theorem}[Fourth-order loss and fewer translates; \stU]\label{spn:thm:source}
Assume \eqref{spn:eq:near} and \eqref{spn:eq:moment} with $u=5b^{-4}$,
$v=9b^{-4}$, and $b\ge6$. Let $S_0\le S$ satisfy
$r_0=s-\dim S_0\le b^{-8}s$. Then some set $A\subseteq G$ satisfies
\begin{align}
 \abs{A}\dim S_0&\le4d\log(1+2b^2),\label{spn:eq:sourcecount}\\
 \dim\sum_{g\in A}C_g(S_0)&\ge\Bigl(1-\frac{20}{b^4}\Bigr)d.\label{spn:eq:sourcerank}
\end{align}
In fact $\abs{A}s\le4d\log(1+2b^2)$.
\end{theorem}

\begin{proof}
Write $L=\log(1+2b^2)$, and choose $t=\floor{4dL/s}$. \Cref{spn:lem:sqrt} gives
$s/d\le1/a$. Since $u\le5/256$ and $v\le9/256$, we have $a\ge251/256$,
$z\le265/256$, and
\[
 0<h=\frac{sz}{2ad}\le\frac{z}{2a^2}\le\frac{33920}{63001}<1 .
\]
Also $4dL/s\ge4aL>1$, so $t\ge1$. In particular $q=1-h$, and
\[
 q^t\le e^{-ht}\le\exp\Bigl(\frac{sz}{2ad}\Bigr)(1+2b^2)^{-2z/a}
 \le\exp\Bigl(\frac{z}{2a^2}\Bigr)(1+2b^2)^{-2z/a}.
\]
The second inequality uses $t\ge4dL/s-1$; this retains the entire cost of the
floor. Applying \cref{spn:lem:restriction} and dividing by $d$, the missing
share is at most
\begin{equation}\label{spn:eq:exactsource}
 \delta_*+\alpha_*\exp\Bigl(\frac{z}{2a^2}\Bigr)(1+2b^2)^{-2z/a}
 +4b^{-8}\log(1+2b^2).
\end{equation}
Bound the three terms separately. First,
$\delta_*=(2u+v-u^2)/(1+v)\le2u+v=19b^{-4}$. Second, $\alpha_*\le1$,
$z/(2a^2)<1$, $e<3$, and $z/a\ge1$, so the middle term is at most
$3(1+2b^2)^{-2}\le\tfrac34b^{-4}$. Third, $\log(1+2b^2)\le2b^2$, so the final
term is at most $8b^{-6}$. For $b\ge6$,
\[
 \frac{19}{b^4}+\frac3{4b^4}+\frac8{b^6}
 \le\frac1{b^4}\Bigl(19+\frac34+\frac8{36}\Bigr)=\frac{719}{36b^4}
 <\frac{20}{b^4}.
\]
This proves \eqref{spn:eq:sourcerank}. Finally,
$\abs{A}\dim S_0\le\abs{A}s\le ts\le4dL$, proving \eqref{spn:eq:sourcecount}.
\end{proof}

\begin{corollary}[Original budget, smaller minimum weight; \stU]\label{spn:cor:five}
For $b\ge4$, the conclusion \eqref{spn:eq:sourcerank} holds with
$\abs{A}s\le5d\log(1+2b^2)$.
\end{corollary}

\begin{proof}
Take $t=\floor{5dL/s}$. The middle term in \eqref{spn:eq:exactsource} becomes
at most $3(1+2b^2)^{-5/2}\le b^{-5}$, and the restriction term is at most
$10b^{-6}$. The total is at most
$b^{-4}(19+1/4+10/16)=159/(8b^4)<20/b^4$.
\end{proof}

\begin{proposition}[The fourth-order scale and its leading coefficient are optimal; \stU]\label{spn:prop:fourth}
For every integer $b\ge2$, the raw error bounds $u=5b^{-4}$, $v=9b^{-4}$, with
$S_0=S$, admit an example whose missing share even after all translates is
\[
 \frac{19b^4-25}{b^4(b^4+9)}=\frac{19}{b^4}-\frac{196}{b^4(b^4+9)}.
\]
Consequently no $o(b^{-4})$ universal deficit bound, and no universal leading
coefficient below $19$, follows from the raw moment and near-isometry
inequalities alone.
\end{proposition}

\begin{proof}
Use \cref{spn:prop:sharp} with $d=b^4(b^4+9)$, $r=(b^4-5)^2$, and $S=\C$. Then
$d-r=19b^4-25$ and $r/d=(1-u)^2/(1+v)$. The unrestricted source has zero
restriction loss. The stated identity follows by division.
\end{proof}

\begin{remark}\label{spn:rem:source-status}
The source's \code{enough_forms} has $(1-4/b^2)d$ in place of
\eqref{spn:eq:sourcerank}, with the count constant five in place of four in
\eqref{spn:eq:sourcecount}. The source imposes much larger lower bounds on $b$
than $b\ge6$, so \cref{spn:thm:source} improves it throughout its admissible
parameter range, using the original moment and ramification estimates and no
new arithmetic input; all three errors, the floor, and the restriction of the
source are retained. The statement has not been executed through Lean. It
improves the tolerance for completing an almost-full span to a basis: with
$O(b^{-4})$ missing dimensions, the per-missing-vector completion cost can be
much larger for a fixed total budget than with $O(b^{-2})$ missing dimensions
(compare \cref{spn:sec:missing}). It does not change the $x^{-1/2+o(1)}$
prime-sum exponent of the source and establishes no further Riemann
hypothesis.
\end{remark}

%======================================================================
\subsection{Exact finite selection certificates}\label{spn:sec:certificates}
\sources{grh_spanning_and_s5_2026-11-26}

The real-valued greedy estimate is convenient for asymptotics. For specified
finite dimensions, integrality gives a more precise computable guarantee, in
rational arithmetic only. It certifies guaranteed ranks, not actual optimizing
translates of an unspecified map.

Let $u,v$ be rational and retain the notation of \cref{spn:sec:spanning}. Put
\[
 R=\ceil*{\frac{(1-u)^2}{1+v}\,d}.
\]
At a current span of dimension $k$, \cref{spn:lem:sqrt} and rank--nullity
guarantee a gain of at least
\[
 g(k)=\max\Bigl\{0,\ s-\floor[\big]{su+s\sqrt{(1+v)k/d}}\Bigr\}.
\]
In particular $g(k)\ge1$ when $k<R$. Define
\[
 T(k)=\min\set{R,k+g(k)},\qquad H(j)=\min_{j\le k\le R}T(k),\qquad
 q_0=0,\quad q_{j+1}=H(q_j).
\]

\begin{proposition}[Envelope certificate; \stU]\label{spn:prop:envelope}
Assume \eqref{spn:eq:near} and \eqref{spn:eq:moment} with rational $u,v$. For
every $j\ge0$ there is a set of at most $j$ translates whose images span a
subspace of dimension at least $q_j$. Moreover $q_{j+1}>q_j$ whenever $q_j<R$,
so $q_j=R$ for some $j\le R$. If the source is restricted to a subspace of
codimension $r_0$, the same $j$ translates retain rank at least $q_j-jr_0$.
\end{proposition}

\begin{proof}
The actual rank can exceed a previous lower certificate, so iterating
$k\mapsto k+g(k)$ on a lower bound would require a monotonicity argument; the
envelope $H$ avoids it. If the actual rank is at least $R$, the target is
already reached. Otherwise the actual integer rank lies between $q_j$ and
$R-1$, so the next actual rank is at least $H(q_j)$. Thus $q_j$ is a lower
certificate after at most $j$ selected translates. Since $g(k)\ge1$ for $k<R$,
we have $T(k)>q_j$ for all $q_j\le k\le R$ when $q_j<R$, which gives strict
growth. The last assertion is \cref{spn:lem:restriction}.
\end{proof}

All square-root comparisons can be eliminated. To decide whether an integer
$m$ is at most $su+s\sqrt{(1+v)k/d}$, set $w=m-su$. The answer is yes if
$w\le0$; if $w>0$, it is yes exactly when $w^2\le s^2(1+v)k/d$. Every quantity
is rational. Binary search over $0\le m\le s$ suffices because values above
$s$ all give gain zero, and the table of $H$ is obtained from one backward pass
over the values of $T$. This constructs the table in $O(R\log(s+1))$ exact
comparisons and $O(R)$ storage.

\begin{example}\label{spn:ex:certificate}
For $d=1000$, $s=20$, and $u=v=1/100$, the sharp universal ceiling is $R=971$.
The exact integer certificate reaches rank $840$ after $148$ translates and
rank $971$ after $261$. The exponential greedy bound in the source, with its
$16\eps$ cutoff and rate $2s/(5d)$, certifies rank $840$ after $230$
translates. These figures compare sufficient guarantees; they do not assert a
lower bound on the optimal number of translates for a particular family. The
certificates, the rational equality examples of \cref{spn:prop:sharp}, the
instantiation at $u=5/b^4$, $v=9/b^4$ and the constants used in
\cref{spn:thm:source} are recorded in a delivered
file.\footnote{\artifact{grh_spanning_and_s5_2026-11-26/checks/spanning_certificates.json}}
\end{example}

The exact least total rank is $r_*=\max\set{\ceil{as},R}$ of
\cref{spn:thm:rank-spectrum}: $R=\ceil{a^2d/z}$ is its second-moment part, and
the other part, $\rank(PB)\ge\ceil{as}$, comes from one map alone. In
\cref{spn:ex:certificate}, $\ceil{as}=20$ and $r_*=R=971$.

%======================================================================

\subsection{Why almost all of the target must be spanned, and how to pay for missing dimensions}\label{src:sec:span}\label{spn:sec:missing}
\begingroup\raggedright
\sources{grh_lean_transfer_analysis; grh_fixed_field_transfer_2026-11-25; grh_transfer_bridges_2026-11-25; grh_beyond_dirichlet_grh_2026-11-25; grh_beyond_dirichlet_2026-11-25; grh_arithmetic_transfers_2026-11-26; grh_transfer_synthesis_2026-11-25; grh_spanning_and_s5_2026-11-26; grh_transfer_advances_2026-11-26; grh_method_transfers_orthogonal_jets_2026-11-26}
\endgroup

The last line of the source comparison divides by $n$ and needs only $n\ge D/2$. The lower
bound (C) needs much more. In \code{integral_forms_long} each vector completing the selected
family to a basis is taken from a globally spanning auxiliary family and has cost
\[
C=2x_1^2b^2+b ,
\]
the first term bounding its squared norm and the second its denominator height. Write
$d=\mathrm{targetDim}$ for the target dimension ($D$ above), $A=\mathrm{mainAmount}$ for the
main amount ($M$ in (A)--(C)) and $u=fa(\log b)^4$, so that $d,C,u$ play the roles of $D,K,U$ in
\cref{src:prop:rank}. Before the rank hypothesis is used, the proof has reached
\begin{equation}\label{spn:eq:source-before-rank}
 0\le c+2\delta+(d-n)C+d(b+u-A),
\end{equation}
which is the second hypothesis of \cref{src:prop:rank} with $r=d-n$, $K=C$ and $B=d(b+u)$.

\paragraph{The source's use.}
Input (B) gives $d-n\le4d/b^2$, and the lemma \code{mul_completion_cost_le} combines the two:
\[
(d-n)C\le4d\bigl(2x_1^2+b^{-1}\bigr)\le db
\]
under the scale inequalities of the source. Since $b\le u$, \eqref{spn:eq:source-before-rank}
yields $c+2\delta\ge d(A-3u)$, which is (C). A fixed positive missing fraction would instead
leave a term of order $dx_1^2b^2$, far above the error budget. Near-full spanning is thus a
requirement of the lower-bound argument even though the final division uses only a positive
fraction of the dimension, and a statement that a fixed fraction of the target is spanned
cannot be substituted for (B) unless the completion is made correspondingly cheap. A sharper
form of (B) is \cref{spn:thm:source}; the weakest completion budget under which the source
argument still runs is \cref{spn:prop:weakening} below.

\begin{example}[Determinants do not control individual directions]\label{src:ex:diag}
The positive forms $\operatorname{diag}(t,t^{-1})$ and the identity on $\R^2$ have the same
determinant, while their restrictions to the first coordinate line have ratio $t$, which is
arbitrary. A bound for a full determinant therefore controls neither the determinant of a
subfamily nor an individual eigenvalue. This is why the lower bound must retain the completing
vectors and their costs, and why the spanning theorem, which controls a large family through
a second moment, does not isolate a prescribed eigenform.
\end{example}

The length $m\asymp\sqrt x$ of the chain is tied to the denominator budget in the same way.
The budget allotted to local defects in the fitting of twists is $fa(\log b)^3/(mk)$, with
$a=xb$ and a further parameter $k$ of the construction, and the explicit estimate for these
defects is $O\bigl(f(m/k)\,b\log b\bigr)$. The quotient of the estimate by the budget is
\begin{equation}\label{src:eq:budget-ratio}
O\Bigl(\frac{m^2}{x(\log b)^2}\Bigr),
\end{equation}
so $m\asymp\sqrt x$ leaves only logarithmic room. A longer chain with the same scale choices
would need better loss estimates, and whenever a normalised loss grows like $m$ the division
by $m$ leaves a constant and the decay disappears. The exponent in \eqref{src:eq:onesided}
cannot be improved by enlarging $m$ with the remaining estimates held fixed.

\paragraph{Paying for missing dimensions.}
By \eqref{src:eq:rank} the deficit is charged by the excess completion cost $K-A$, not by $K$
alone. Normalizing all errors by one scale gives the following form.

\begin{corollary}[Completion cost at a common error scale; \stU]\label{spn:cor:rank-error}
In \cref{src:prop:rank} write $\varrho=n/D$ and let $E_0>0$. Suppose
\[
 \varrho\ge\rho>0,\quad U\le U_0E_0,\quad \delta/D\le I_0E_0,\quad
 B/D\le J_0E_0,\quad(1-\varrho)(K-A)_+\le K_0E_0 .
\]
Then
\begin{equation}\label{spn:eq:rank-error}
 R\le\frac{U_0+(2I_0+J_0+K_0)/\rho}{m}\,E_0 .
\end{equation}
\end{corollary}

\begin{proof}
Divide \eqref{src:eq:rank} by $nm$ and use $D/n=1/\varrho\le1/\rho$ and
$r/n=(1-\varrho)/\varrho$.
\end{proof}

There are consequently two ways to achieve the needed accuracy: make the
missing dimension fraction extremely small, or prove that a completion has
excess cost $O(E_0)$. In the latter case a fixed positive fraction of
independent target forms can suffice. If $K\le A$, the deficit term is
favorable. A uniform lower bound on $n/D$, or a compensating improvement in
the other errors, is still needed; the resulting exponent bookkeeping is
\cref{src:rem:rates}.

\begin{proposition}[An explicit weakening inside the source proof; \stU]\label{spn:prop:weakening}
In the notation of \eqref{spn:eq:source-before-rank}, the source's conclusion
$c+2\delta\ge d(A-3u)$ follows from \eqref{spn:eq:source-before-rank} under the completion
budget
\begin{equation}\label{spn:eq:weaker-spanning}
 (d-n)C\le d(2u-b).
\end{equation}
Keeping the same completion and height estimates, it is sufficient to replace
the loss bound $d-n\le4d/b^2$ by
\[
 \frac{d-n}{d}\le\min\Bigl\{\frac12,\ \frac{u}{2x_1^2b^2+b}\Bigr\},
\]
and with $A\ge0$, $\delta\le dI_0$, $b\le u$ and the short-family upper bound
$c\le n(A-mR+u)$ of \cref{src:prop:rank} this gives $mR\le7u+4I_0$.
\end{proposition}

\begin{proof}
Inequality \eqref{spn:eq:source-before-rank} gives
$c+2\delta\ge d(A-b-u)-(d-n)C$, which is at least $d(A-3u)$ exactly when
\eqref{spn:eq:weaker-spanning} holds. Since the source has $b\le u$, the
condition $(d-n)C\le du$ suffices; this is the second term of the displayed
minimum. The first term ensures $n\ge d/2$. For the final bound use
$B=d(b+u)$ in \eqref{src:eq:rank}, drop the favorable term $-(d-n)A$, and
bound $2\delta+B+(d-n)C$ by $d(2I_0+3u)$.
\end{proof}

The completing vectors of the source come from a globally spanning auxiliary
family, with squared-norm exponent $2x_1^2b^2$, scalar denominator-height
bound $b$, and the displayed volume-height bound; an arbitrary completion is
insufficient. Compared with the allowance $(d-n)C\le db$ of the source's use,
\cref{spn:prop:weakening} is a weakening of a sufficient hypothesis in the
written source argument, not an improvement to the endpoint
\cref{src:thm:dgrh}: for the same expensive completion the missing fraction
still tends to zero. For a new target it isolates a distinct task: construct
cheaper integral completing vectors and control their twist scalars, instead
of reproducing the full spanning estimate. The abstract inequality still
requires injective twist maps, a family and a completion that together form a
basis, integral vectors after a common scalar, norm bounds at all complex
conjugates, and height bounds at the allowed finite places, all tied to the
same lattice. With the source-like error scale $E_0=fa(\log b)^4$ and
$m\asymp\sqrt x$, \eqref{spn:eq:rank-error} has normalized size
$O((\log b)^4/\sqrt x)$; converting this into a prime remainder still needs
the comparison of the arithmetic gain with the relevant Frobenius statistic.

%======================================================================
\subsection{Optimal integral frames and the ramification obstruction}\label{spn:sec:frames}
\sources{grh_nonabelian_transfer_2026-11-25}

For a totally real number field $F$ of degree $n$, an \emph{ordinary integral
tight frame} of bound $M\in\Z_{>0}$ is a finite family $v_i\in\OO_F$ with
positive integer multiplicities $m_i$ such that
\begin{equation}\label{spn:eq:frame}
 \sum_im_i\sigma(v_i)\tau(v_i)=M\delta_{\sigma,\tau}
 \qquad(\sigma,\tau\colon F\hookrightarrow\R).
\end{equation}
Thus the vectors give a resolution of a scalar identity in the embedding
space. Repeated vectors can replace multiplicities. This is the linear-algebra
part of the source's frame input; its global integrality condition instead
uses a fractional Weil lattice (\cref{spn:sec:projective}).

\begin{theorem}[Trace-dual obstruction; \stU]\label{spn:thm:dual-obstruction}
If \eqref{spn:eq:frame} holds, then
\begin{equation}\label{spn:eq:codiff-obstruction}
 M\,\spnDD_{F/\Q}^{-1}\subseteq\OO_F .
\end{equation}
In particular, every rational prime ramified in $F$ divides $M$.
\end{theorem}

\begin{proof}
Let $a$ lie in the codifferent, so $\Tr_{F/\Q}(av_i)\in\Z$ for all $i$.
Applying the frame operator to the embedding vector of $a$ gives
\[
 Ma=\sum_im_i\Tr_{F/\Q}(av_i)v_i\in\OO_F .
\]
This proves \eqref{spn:eq:codiff-obstruction}. At a prime $\mathfrak p$ over a
ramified rational prime $p$, write $e_{\mathfrak p}$ for the ramification
index and $d_{\mathfrak p}>0$ for the different exponent. The ideal
containment requires $e_{\mathfrak p}\vp(M)\ge d_{\mathfrak p}$, so
$\vp(M)\ge1$.
\end{proof}

The obstruction concerns every ordinary integral frame, not only a particular
construction. It also explains why the source allows denominators in one half
of the relative codifferent: those denominators change the dual lattice. The
elementary obstruction has a converse away from a finite initial range.

\begin{theorem}[Eventual frame bounds for an integral lattice; \stU]\label{spn:thm:real-eventual}
Let $L$ be a positive definite integral $\Z$-lattice of positive rank with
Gram matrix $G$ in an integral basis. Put
\[
 d=\min\set{a\in\Z_{>0}: aG^{-1}\text{ has integer entries}} .
\]
Every positive integer tight-frame bound for vectors in $L$ is divisible by
$d$, and every sufficiently large positive multiple of $d$ is such a bound.
\end{theorem}

\begin{proof}
Let $B$ be the real basis matrix, so $G=B^{\spnT}B$. If the frame vectors have
integer coordinate columns $u_i$, the frame identity becomes
\[
 \sum_im_iu_iu_i^{\spnT}=MG^{-1}.
\]
The left side is integral. This proves $d\mid M$; equivalently, $d$ is the
exponent of the discriminant group $L^\#/L$.

For sufficiency set $C=dG^{-1}$, a positive definite integral symmetric
matrix. The rank-one matrices generated by $e_i$ and $e_i+e_j$ span all
integral symmetric matrices over $\Z$: their differences give the
off-diagonal matrix units, while $e_ie_i^{\spnT}$ gives the diagonal units.
Let $S$ be this finite set of vectors. For a sufficiently small positive
rational $\epsilon$, the matrix
\[
 C-\epsilon\sum_{v\in S}vv^{\spnT}
\]
is positive definite and rational. Rational $LDL^{\spnT}$ decomposition
expresses it as a sum of positive rational multiples of rational rank-one
matrices. Scaling their vectors to integer vectors and clearing coefficients
yields
\[
 qC=\sum_{v\in T}c_vvv^{\spnT},\qquad c_v\in\Z_{>0},\quad S\subseteq T .
\]
Write $C=\sum_{v\in T}z_vvv^{\spnT}$ with $z_v\in\Z$, using $S$, and extending
by zero on other vectors. For $N=kq+r$, $0\le r<q$,
\[
 NC=\sum_{v\in T}(kc_v+rz_v)vv^{\spnT}.
\]
For all sufficiently large $k$, every coefficient is nonnegative. The
resulting frame has bound $Nd$. The same argument makes every coefficient
exceed any prescribed fixed lower bound.
\end{proof}

This proof supplies an effective threshold once a rational reserve
decomposition has been chosen. It supplies no useful uniform bound as the rank
and arithmetic data vary. Sums of integral squares and Hermitian norms form an
established subject \cite{Beli2019Waring}; the present formulation is used as a
transfer tool.

\begin{example}[The cubic field of discriminant $229$]\label{spn:ex:cubic229}
Let $f(X)=X^3-4X+1$ and let $\alpha$ be a root. The rational root test proves
irreducibility. There is a root in each of $(-3,-2)$, $(0,1)$ and $(1,2)$, so
the cubic field $F_0=\Q(\alpha)$ is totally real. Its polynomial discriminant
is $-4(-4)^3-27=229$, a prime. Thus $(1,\alpha,\alpha^2)$ is an integral
basis. Newton sums give
$\Tr(1),\Tr(\alpha),\dots,\Tr(\alpha^4)=(3,0,8,-3,32)$, and the trace Gram
matrix is
\[
 T=\begin{pmatrix}3&0&8\\0&8&-3\\8&-3&32\end{pmatrix},\qquad
 229\,T^{-1}=\begin{pmatrix}247&-24&-64\\-24&32&9\\-64&9&24\end{pmatrix}.
\]
The coordinate vectors
\[
 (0,2,1),\quad(1,-1,1)\ \text{twice},\quad(2,-1,1),\quad(4,-5,-2),\quad
 (15,0,-4)
\]
in the basis $(1,\alpha,\alpha^2)$ have outer products summing to
$229\,T^{-1}$. Therefore they give a six-vector frame of bound $229$, and
\cref{spn:thm:dual-obstruction} shows that no smaller positive integer bound
is possible.
\end{example}

Let now $F$ be the splitting field of $f$ and let $\alpha,\beta$ be distinct
roots. Its Galois group is $S_3$, because the irreducible cubic has nonsquare
discriminant. Its six embeddings correspond to ordered pairs of distinct
roots. Define
\[
 b=(1,\alpha,\alpha^2,\beta,\alpha\beta,\alpha^2\beta).
\]
This is an integral basis. Indeed, only $229$ can ramify in the splitting
field. At $229$ the extension is tame, and
\[
 f(X)\equiv(X-29)^2(X+58)\pmod{229}.
\]
Inertia has order two and acts as a transposition: it is nontrivial, cyclic,
and fixes the simple-root branch. A tame degree-six Galois extension with
inertia index two has discriminant exponent $6(1-1/2)=3$. Consequently
$\disc(F)=229^3$. The Gram matrix of $b$, computed below, also has determinant
$229^3$, proving that the lattice it generates has index one in $\OO_F$.

Write $s_j$ for the $j$th power sum of the three roots. Then $s_0=3$, $s_1=0$,
$s_2=8$ and $s_j=4s_{j-2}-s_{j-3}$ for $j\ge3$. For monomials
$\alpha^i\beta^j$ and $\alpha^k\beta^l$, the trace of their product equals
$s_{i+k}s_{j+l}-s_{i+j+k+l}$. It follows that the Gram matrix and its scaled
inverse are
\begin{align*}
G&=\begin{pmatrix}
6&0&16&0&-8&3\\
0&16&-6&-8&3&-32\\
16&-6&64&3&-32&20\\
0&-8&3&16&3&32\\
-8&3&-32&3&32&-4\\
3&-32&20&32&-4&125
\end{pmatrix},\\[3pt]
C:=229\,G^{-1}&=\begin{pmatrix}
125&-16&-32&4&0&-3\\
-16&32&3&0&-3&8\\
-32&3&16&-3&8&0\\
4&0&-3&32&-6&-8\\
0&-3&8&-6&16&0\\
-3&8&0&-8&0&6
\end{pmatrix}.
\end{align*}

\begin{theorem}[Optimal discriminant-$229$ frame; \stU]\label{spn:thm:optimal229}
Let $v_i$ have the following coordinates in $b$, with multiplicities $m_i$:
\[
\begin{array}{@{}rrrrrrcr@{}}
\toprule
1&\alpha&\alpha^2&\beta&\alpha\beta&\alpha^2\beta&&m_i\\\midrule
0&2&-2&-2&-1&1&&1\\
1&-2&-2&4&-2&-1&&1\\
2&2&0&1&2&1&&1\\
4&1&-1&1&-1&0&&2\\
4&1&-1&2&0&-1&&1\\
6&-4&-1&-1&2&-1&&1\\
6&-1&-2&-2&-1&1&&1\\\bottomrule
\end{array}
\]
They satisfy \eqref{spn:eq:frame} with $M=229$, and $229$ is the smallest
possible positive integer bound for any ordinary integral frame in $F$.
\end{theorem}

\begin{proof}
Let $u_i$ denote the displayed columns. Direct integer multiplication gives
\[
 \sum_im_iu_iu_i^{\spnT}=C,\qquad GC=229\,I_6 .
\]
If $B=(\sigma(b_j))_{\sigma,j}$, then $B^{\spnT}B=G$, whence
\[
 \sum_im_i\sigma(v_i)\tau(v_i)=(BCB^{\spnT})_{\sigma,\tau}
 =229\,\delta_{\sigma,\tau}.
\]
The vectors are integral. Every positive integer frame bound is divisible by
the ramified prime $229$, by \cref{spn:thm:dual-obstruction}. This proves
optimality.
\end{proof}

The squared trace norms of the seven displayed vectors are, in order,
\[
 (221,\ 229,\ 229,\ 74,\ 157,\ 211,\ 179),
\]
and their multiplicity-weighted sum is $6\cdot229$. In particular all real
conjugates of every displayed vector have absolute value at most
$\sqrt{229}$. The Newton sums, matrices, determinants, outer products and
identities are recomputed in an exact
certificate.\footnote{\artifact{grh_nonabelian_transfer_2026-11-25/certificates/frame_certificates.json}}

\begin{corollary}[Cyclotomic tensor frame; \stU]\label{spn:cor:tensor}
For $229\mid q$, the elements $v_i\zeta_q^j$, $0\le j<q$, with the same
multiplicities, give a frame of bound $229q$ on the compatible embeddings of
$F\Q(\zeta_q)$. The same embedding identity holds on the induced CM type after
adjoining the source's imaginary quadratic field $K$.
\end{corollary}

\begin{proof}
Multiply \eqref{spn:eq:frame} by the geometric-sum identity
\[
 \sum_{j=0}^{q-1}\eta(\zeta_q^j)\overline{\eta'(\zeta_q^j)}
 =q\,\delta_{\eta,\eta'} .
\]
Embeddings of the compositum correspond to compatible pairs of embeddings.
Restricting the product identity to those pairs proves the assertion; linear
disjointness is not assumed.
\end{proof}

The tensor frame preserves integrality at the inert primes used in the local
argument for the nonabelian benchmark of \cref{grp:sec:s3}. Its bound still
carries $229$. The correction needed for coprime bounds is a change to the
appropriate fractional lattice, addressed in \cref{spn:sec:projective}.

%======================================================================

\subsection{Integral Gram determinants with vanishing minors}\label{src:sec:gram}
\begingroup\raggedright
\sources{grh_new_results_2026-11-26; gram-euler-and-zero-filters; gram_column_number_2026-11-27; proveit_three_extensions_2026-11-27_1}
\endgroup

Input (C) is a lower bound for the Gram determinant of an integral family, and inputs (A) and
(B) ask for many short independent forms. A finite-dimensional question isolates the tension
between these requirements: how small can the Gram determinant of an integer matrix be when
many of its maximal minors are nonzero? Cauchy--Binet and integrality give one lower bound; the
Pl\"ucker relations between the minors force a stronger one, with an explicit penalty for
vanishing minors. The question complements the integral frames of \cref{spn:sec:frames}: there
the frame operator of an integral family is required to be a scalar, and the trace-dual
obstruction (\cref{spn:thm:dual-obstruction}) and \cref{spn:thm:real-eventual} determine which
scalars occur; here only the Gram determinant is controlled, and the question is how small it
can be.

\begin{theorem}[Arithmetic determinant amplification; \stU]\label{src:thm:gram}
Let $A\in\Z^{r\times n}$ have rank $r$, where $n\ge5$ and $2\le r\le n-2$. For an $r$-element
set $B\subseteq\set{1,\dots,n}$ let $A_B$ be the square matrix of the columns in $B$, in
increasing order, and put
\[
 N=\binom nr,\qquad z=\#\set{B:\det A_B=0},\qquad
 D=\gcd\set{\abs{\det A_B}:\det A_B\ne0}.
\]
Then
\begin{equation}\label{src:eq:gram}
 \frac{\det(AA^{\srctr})}{D^2}\ \ge\ \max\Bigl\{N-z,\ \frac85N-4z,\ \frac{12}5N-12z\Bigr\}.
\end{equation}
In particular, if every $r$ columns are linearly independent,
$\det(AA^{\srctr})\ge\tfrac{12}5\binom nrD^2$. The coefficient $12/5$ cannot be increased
uniformly over this class, and each of the three affine bounds in \eqref{src:eq:gram} is
attained.
\end{theorem}

The theorem does not assert that the maximum in \eqref{src:eq:gram} is optimal for every fixed
triple $(r,n,z)$. Cauchy--Binet and integrality alone give only $N-z$. On four labels the
relevant relation is
\begin{equation}\label{src:eq:plucker}
 p_{12}p_{34}-p_{13}p_{24}+p_{14}p_{23}=0,
\end{equation}
which holds for the two-row minors of any matrix by direct expansion.

\begin{lemma}[Local Pl\"ucker bounds; \stU]\label{src:lem:plucker}
\begin{enumerate}[label=\textup{(\roman*)}]
\item If the six integers in \eqref{src:eq:plucker} are nonzero, their squares sum to at least
$9$.
\item If ten nonzero integers $p_{ij}$, $1\le i<j\le5$, satisfy \eqref{src:eq:plucker} on every
four-element set of labels, their squares sum to at least $24$.
\item For ten integers $p_{ij}$, $1\le i<j\le5$, satisfying \eqref{src:eq:plucker} on every
four labels, let $w$ be the number of zeros and $V$ the sum of squares. Then
$V\ge\max\set{10-w,\ 16-4w,\ 24-12w}$.
\end{enumerate}
\end{lemma}

\begin{proof}
(i) If all six absolute values were one, the left side of \eqref{src:eq:plucker} would be a
signed sum of three ones and could not vanish; so one square is at least four.

(ii) Suppose the sum is at most $23$. No coordinate has absolute value four or more, since
that contributes at least $16+9=25$, and at most one has absolute value three, since two
contribute at least $18+8=26$. If the edge $\set{a,b}$ has absolute value three, at most one
edge has absolute value two, because relative to the baseline $10$ the three adds $8$ and two
twos would add $6$. Choose four labels containing $\set{a,b}$ and avoiding that edge, by
omitting one of its endpoints outside $\set{a,b}$; the three absolute products in the relation
are then $3,1,1$, which have no vanishing signed sum. So all absolute values are one or two.
Call an edge heavy if its absolute value is two. On any four labels the absolute products lie
in $\set{1,2,4}$, and a signed sum of three nonzero reals vanishes only if one absolute value is
the sum of the other two, so the product multiset is $\set{2,1,1}$ or $\set{4,2,2}$, with one or
four heavy edges among the six. There are five four-label subsets and each edge lies in three of
them. If $h$ edges are heavy and $k$ subsets have four heavy edges, double counting gives
$3h=(5-k)+4k=5+3k$, which is impossible modulo three.

(iii) Integrality gives $V\ge10-w$, and (ii) gives $V\ge24$ for $w=0$. If $w=1$, delete an
endpoint of the zero edge: the six coordinates on the remaining four labels are nonzero and
contribute at least $9$ by (i), and the three other nonzero coordinates at least $3$, so
$V\ge12$. These cases prove $V\ge24-12w$, whose right side is nonpositive for $w\ge2$, and
$V\ge16-4w$ for $w\le1$; for $w\ge2$, $10-w\ge16-4w$.
\end{proof}

\begin{proof}[Proof of \cref{src:thm:gram}]
Put $p_B=\det(A_B)/D\in\Z$. By Cauchy--Binet,
$S:=\det(AA^{\srctr})/D^2=\sum_{\abs B=r}p_B^2$. First let $r\le n-3$. For disjoint sets
$I,J$ with $\abs I=r-2$ and $\abs J=5$, and distinct $a,b\in J$, let
$q_{ab}=D^{-1}\det(A_{I,a,b})$, the columns of $I$ in increasing order followed by $a,b$. These
are signed versions of $p_{I\cup\set{a,b}}$ and satisfy the four-label relations: if the
columns of $I$ are dependent all $q_{ab}$ vanish, and otherwise an invertible real row
operation sends them to the first $r-2$ standard basis vectors, and expansion along those
columns writes the $q_{ab}$ as a common real multiple of two-row minors. Let $V_{I,J}$ and
$w_{I,J}$ be the local sum of squares and number of zeros. There are
$\binom n{r-2}\binom{n-r+2}5$ local systems, and each $p_B$ occurs in exactly
$t=\binom r2\binom{n-r}3>0$ of them: choose the two elements of $B$ outside $I$, then the three
elements of $J$ outside $B$. Hence $10\binom n{r-2}\binom{n-r+2}5=Nt$,
$\sum_{I,J}V_{I,J}=tS$ and $\sum_{I,J}w_{I,J}=tz$. Summing each of the three affine
inequalities of \cref{src:lem:plucker}(iii) and dividing by $t$ gives the three bounds.

If $r=n-2$, the complementary coordinates $p^*_{ij}=\pm p_{\set{1,\dots,n}\setminus\set{i,j}}$,
with the signs of the complementary permutation, are, up to one common nonzero real factor, the
two-row minors of a basis of the two-dimensional orthogonal complement of the row space of $A$,
by Laplace expansion. So they are integers satisfying the rank-two relations, with the same
sum of squares and number of zeros, and $\binom n2=N$; apply the case $r=2$. The examples below
give the sharpness assertions.
\end{proof}

\begin{example}[Equality cases]\label{src:ex:gram}
Let $e_1=(1,0)^{\srctr}$, $e_2=(0,1)^{\srctr}$, $u=(1,1)^{\srctr}$, $v=(1,-1)^{\srctr}$
and $w=(1,2)^{\srctr}$. The following matrices have $r=2$, $n=5$, $N=10$ and $D=1$.
\begin{center}\small
\begin{tabular}{cllc}
\toprule
$z$ & columns of $A$ & $AA^{\srctr}$ & $\det(AA^{\srctr})$\\
\midrule
$0$ & $e_1,e_2,u,v,w$ & $\left(\begin{smallmatrix}4&2\\2&7\end{smallmatrix}\right)$ & $24$\\
$1$ & $e_1,e_1,e_2,u,v$ & $\left(\begin{smallmatrix}4&0\\0&3\end{smallmatrix}\right)$ & $12$\\
$2$ & $e_1,e_1,e_2,e_2,u$ & $\left(\begin{smallmatrix}3&1\\1&3\end{smallmatrix}\right)$ & $8$\\
$3$ & $e_1,e_1,e_1,e_2,u$ & $\left(\begin{smallmatrix}4&1\\1&2\end{smallmatrix}\right)$ & $7$\\
\bottomrule
\end{tabular}
\end{center}
The first attains $12N/5-12z$, the second both nontrivial bounds, the third $8N/5-4z$ and the
baseline, and the fourth the baseline alone. The rows of
\[
 \begin{pmatrix}-1&-1&1&0&0\\-1&1&0&1&0\\-1&-2&0&0&1\end{pmatrix}
\]
form an integral basis of the kernel of the first example; this rank-three matrix has ten
nonzero maximal minors with gcd one and Gram determinant $24$. Exact integer arithmetic confirms
all five values, Cauchy--Binet and every maximal minor
(\artifact{grh_new_results_2026-11-26/artifacts/determinant_results.json}).
\end{example}

With $\delta=z/N$ the normalized bound is $12/5-12\delta$ for $\delta\le1/10$,
$8/5-4\delta$ for $1/10\le\delta\le1/5$, and $1-\delta$ for $\delta\ge1/5$.

\begin{corollary}[Short integral columns are limited; \stU]\label{src:cor:gram}
Under the hypotheses of \cref{src:thm:gram} let $E=\sum_{j=1}^n\norm{A_j}^2$ be the total
squared length of the columns and $F$ the right side of \eqref{src:eq:gram}. Then
$E\ge rD^{2/r}F^{1/r}$. If every maximal minor is nonzero and every column has length at most
$H$, then $(nH^2/r)^r\ge\tfrac{12}5\binom nrD^2$.
\end{corollary}

\begin{proof}
The eigenvalues of $AA^{\srctr}$ are positive with sum $E$, so the arithmetic--geometric mean
inequality gives $(E/r)^r\ge\det(AA^{\srctr})$; apply \cref{src:thm:gram}, and $E\le nH^2$.
\end{proof}

A common determinantal divisor, many independent column subsets and a small total squared
length are therefore competing requirements. This does not by itself improve the comparison of
\cref{src:sec:comparison}: an application would have to match the forms and lattices of the
source to the hypotheses and the normalization used here. The weaker fact that a generic
two-row integer matrix with five columns has a minor of absolute value at least three is known
from the study of generic integer matrices with bounded maximal minors
\cite{KriepkeKyureghyanSchymura2023Size,KriepkeSchymura2025Generic}; the aggregate
sum-of-squares inequality and its dependence on vanishing minors are the content of
\cref{src:thm:gram}.

\paragraph{The rank-two minimum for many columns.} For $r=2$ and $z=0$ the least value is known up to
$O(n^2)$. For $A=[v_1\ \cdots\ v_n]\in\Z^{2\times n}$ with $\det(v_i,v_j)\ne0$ for $i\ne j$ put
\[
 D(A)=\gcd_{i<j}\abs{\det(v_i,v_j)},\qquad E(n)=\min_A\frac{\det(AA^{\srctr})}{D(A)^2},
\]
so that \cref{src:thm:gram} gives $E(n)\ge\frac{12}5\binom n2$ for $n\ge5$, with equality at $n=5$. Let
$w_j=\varphi(j)/j$ and
\begin{equation}\label{src:eq:gramPhi}
 \Phi(T)=\frac4{3T}\sum_{1\le j\le T}w_j\Bigl(1-\frac{j^2}{T^2}\Bigr)^{3/2}\qquad(T>0),
\end{equation}
the empty sum being zero, and $\Phi_*=\max_{T>0}\Phi(T)$.

\begin{theorem}[Asymptotic rank-two minimum; \stU]\label{src:thm:gramasym}
The maximum $\Phi_*$ is attained at a unique $T_*\in(6,7)$. With $\kappa=(64\Phi_*^2)^{-1}$,
\begin{equation}\label{src:eq:gramasym}
 E(n)=\kappa(n-1)^4+O(n^2),\qquad0.068157932354<\kappa<0.068157932355 .
\end{equation}
For all sufficiently large $n$ and every minimizing matrix there is a basis of the lattice generated by its
columns in which, after independent sign changes, the columns are $(1,0)^{\srctr}$ and primitive vectors
$(m,j)^{\srctr}$ on exactly the rows $1\le j\le6$, the number of columns on row $j$ being
\begin{equation}\label{src:eq:gramrows}
 N_j=\frac{w_j}{\Phi_*T_*}\sqrt{1-j^2/T_*^2}\,(n-1)+O(1).
\end{equation}
All implied constants are absolute.
\end{theorem}

The weights in \eqref{src:eq:gramrows} are approximately $0.32484$, $0.15610$, $0.19327$, $0.12773$,
$0.16219$, $0.03587$. B\'ar\'any and Tokushige study the sum of the absolute pairwise determinants in the
problem of the minimum area of a convex lattice polygon \cite{BaranyTokushige2004Minimum}, and Kriepke and
Schymura the largest absolute minor \cite{KriepkeSchymura2025Generic}; for the sum of squares a trace
minimization gives both the scalar optimization and the arithmetic profile.

\emph{Normalization and change of metric.} Let $\Lambda_A=\sum_i\Z v_i$. Its index in $\Z^2$ is $D(A)$, by
the Smith normal form. In a basis of $\Lambda_A$ the columns generate $\Z^2$ and every minor is divided by
$D(A)$ up to a common sign, and Cauchy--Binet gives $\det(AA^{\srctr})=\sum_{i<j}\det(v_i,v_j)^2$. So $E(n)$
is also the least value of the unnormalized sum $\sum_{i<j}\det(v_i,v_j)^2$ over $n$ pairwise nonparallel
integer vectors; the minimum exists, and each column of a minimizing configuration that generates $\Z^2$ is
primitive, since dividing a vector by its content strictly decreases the sum. For a positive definite
$2\times2$ matrix $M$,
\begin{equation}\label{src:eq:gramtrace}
 \min_{g\in\SL_2(\R)}\tr(gMg^{\srctr})=2\sqrt{\det M},
\end{equation}
by the arithmetic--geometric mean inequality, with equality at $g=(\det M)^{1/4}M^{-1/2}$. For a lattice
$\Lambda\subset\R^2$ of covolume one let $\mathcal P(\Lambda)$ contain one vector from each antipodal pair
of primitive vectors. Applying \eqref{src:eq:gramtrace} to $M=\sum_iv_iv_i^{\srctr}$ gives
\begin{equation}\label{src:eq:grammin}
 2\sqrt{E(n)}=\inf_\Lambda\ \min_{S\subset\mathcal P(\Lambda),\ \abs S=n}\ \sum_{v\in S}\abs v^2 .
\end{equation}
Any $n\ge2$ distinct directions span; in a fixed lattice the inner minimum takes $n$ shortest primitive
directions, and an isotropized minimizing matrix realizes the joint infimum.

\emph{Counting primitive directions.} Put
$P_\Lambda(R)=\sum_{v\in\mathcal P(\Lambda)}(R^2-\abs v^2)_+$ and
$N_\Lambda(R)=\#\set{v\in\mathcal P(\Lambda):\abs v\le R}$. After a rotation a reduced basis of $\Lambda$ is
$(a,0),(a\theta,1/a)$, with $a$ the least nonzero length and $\abs\theta\le1/2$; the first vector extends to a
basis, and reducing the second modulo the first gives $a^2\le a^2/4+a^{-2}$, so $a\le(4/3)^{1/4}$. Taking
the representative with positive second coordinate on each nonhorizontal pair, row $j>0$ consists of the
vectors $(a(m+j\theta),j/a)$ with $\gcd(m,j)=1$, and row zero has one representative. Put $T=aR$.

\begin{lemma}[Truncated sums of primitive directions; \stU]\label{src:lem:gramcount}
Uniformly in $\Lambda$, as $R\to\infty$,
\begin{equation}\label{src:eq:gramcrude}
 P_\Lambda(R)=R^4\Phi(T)+O\bigl(R^3\log(2R)\bigr).
\end{equation}
For every compact interval $I\subset(0,\infty)$, uniformly for $T\in I$,
\begin{equation}\label{src:eq:gramsharp}
 P_\Lambda(R)=R^4\Phi(T)+R^2+O_I(1),\qquad N_\Lambda(R)=R^2\beta(T)+1+O_I(1),
\end{equation}
where $\beta(T)=\frac2T\sum_{1\le j\le T}w_j\sqrt{1-j^2/T^2}=2\Phi(T)+\frac T2\Phi'(T)$.
\end{lemma}

\begin{proof}
For $1\le j\le T$ put $\eta=(R^2-j^2/a^2)^{1/2}$. In the horizontal coordinate the summand on row $j$ is
$f(x)=(\eta^2-x^2)_+$, whose integral is $4\eta^3/3$. The mean density of the coprime indices on the row is
$w_j/a$, so the integral terms sum to $R^4\Phi(T)$.

For the uniform error write $\mathbf1_{\gcd(m,j)=1}=\sum_{d\mid\gcd(m,j)}\mu(d)$. On a shifted grid of
spacing $ad$ the sum of $f$ differs from its integral divided by $ad$ by at most the total variation
$2\eta^2$, so the error on row $j$ is $O(R^2\tau(j))$, where $\tau$ counts divisors. Since
$\sum_{j\le T}\tau(j)\le T(1+\log T)$ for $T\ge1$ and $T=O(R)$, summing and adding the horizontal term
proves \eqref{src:eq:gramcrude}, also for $T<1$.

For the finer error, a shifted grid of spacing $d_0$ satisfies
$\abs{\sum_kf(d_0k+u)-d_0^{-1}\int f}\le d_0\operatorname{Var}(f')$: integrate the linear interpolant on
each cell, whose slope is the average of $f'$ there, so the error on a cell is at most $d_0^2$ times the
variation of $f'$ on it. Here $\operatorname{Var}(f')=8\eta$. Split row $j$ into its $\varphi(j)$ primitive
residue classes modulo $j$ and take $d_0=aj$. Only boundedly many rows and classes occur, since $aR=T$ is
bounded, so the total error on positive rows is $O_I(1)$, and the horizontal point contributes
$R^2-a^2=R^2+O_I(1)$. Counting the points of each residue grid in $[-\eta,\eta]$ as its length divided by
$aj$, with error $O(1)$, gives the count. The identity for $\beta$ follows by differentiation, at integers
by continuity.
\end{proof}

\begin{lemma}[The certified scalar maximum; \stU]\label{src:lem:grammax}
The maximum of $\Phi$ is attained at the unique $T_*\in[6.15,6.60]$ with
\begin{equation}\label{src:eq:gramroot}
 \sum_{j=1}^6w_j(4j^2-T_*^2)\sqrt{T_*^2-j^2}=0 ,
\end{equation}
it is nondegenerate, $\Phi(T)<0.4787$ outside $[6.15,6.60]$, and
\[
 6.349265427728<T_*<6.349265427729,\qquad0.478797428690<\Phi_*<0.478797428691 .
\]
\end{lemma}

\begin{proof}
Let $A_0(x)=\sum_{j\le x}w_j=\sum_{d\le x}\mu(d)\floor{x/d}/d$. Comparison with $6x/\pi^2$, bounding the
floor errors by the harmonic sum and the tail by $x/\floor x$, gives $\abs{A_0(x)-6x/\pi^2}\le\log x+3$
for $x\ge1$. Partial summation against $(1-x^2/T^2)^{3/2}$, whose integral over $[0,T]$ is $3\pi T/16$ and
whose total decrease is one, gives
$\abs{\Phi(T)-3/(2\pi)}\le4(\log T+3)/(3T)$ for $T\ge1$, which is below $0.4787-3/(2\pi)$ for $T\ge20000$.
Two finite checks in integer arithmetic close the remaining ranges. First, the totient sieve with upward
rounding of $10^9w_j$ certifies $A_0(j)\le76j/125+3/5$ for $1\le j\le20000$, hence between integers by
monotonicity of the right side, and partial summation gives
$\Phi(T)\le(76/125)\pi/4+4/(5T)<0.4787$ for $1000\le T\le20000$. Second, $[1,1000]\setminus(6.15,6.60)$ is
subdivided into intervals $[l,r]$; the summands of $T\Phi(T)$ are nondecreasing in $T$, so
$\Phi(T)\le r\Phi(r)/l$, and with every radical rounded upwards by integer square roots and bisection of
every interval whose bound is at least $0.4787$, all $3255$ terminal intervals pass. Below $1$ the sum is
empty.

A lower evaluation gives $\Phi(6.35)>0.47879$. With $u=T^2$ the derivative of $\Phi$ on $(6,7)$ has the sign
of $\Gamma(u)=\sum_{j=1}^6w_j(4j^2-u)\sqrt{u-j^2}$, and
$\Gamma''(u)=-\frac{3u}4\sum_{j=1}^6w_j(u-j^2)^{-3/2}<0$. The signs of $\Gamma$ at $u=6.15^2$ and $6.60^2$ are
positive and negative, certified by the same rational bounds on the radicals. Strict concavity gives a
unique zero with $\Gamma'<0$ there, hence $\Phi''(T_*)<0$, and bisection with the monotone bound gives the
intervals for $T_*$, $\Phi_*$ and $\kappa$.
\end{proof}

{\raggedright The finite checks are carried out by
\artifact{gram-euler-and-zero-filters/gram/verify_constant.py}, with output
\artifact{gram-euler-and-zero-filters/gram/constant_verification.json}.\par}

\begin{proof}[Proof of \cref{src:thm:gramasym}]
For every set $S$ of $n$ directions in $\Lambda$ and every $R>0$,
\begin{equation}\label{src:eq:gramduality}
 \sum_{v\in S}\abs v^2\ge nR^2-P_\Lambda(R).
\end{equation}
Take $R^2=(n-1)/(2\Phi_*)$ and $I=[6.15,6.60]$. For $T\in I$, \eqref{src:eq:gramsharp} gives
$P_\Lambda(R)\le\Phi_*R^4+R^2+O(1)$; for $T\notin I$ the error $o(R^4)$ in \eqref{src:eq:gramcrude} is
absorbed, for large $R$, by the gap $\Phi_*-0.4787>0$ of \cref{src:lem:grammax}. With
\eqref{src:eq:grammin},
\[
 2\sqrt{E(n)}\ge(n-1)R^2-\Phi_*R^4-O(1)=\frac{(n-1)^2}{4\Phi_*}-O(1).
\]
For the upper bound take $a=T_*/R$ and the rectangular lattice with basis $(a,0),(0,1/a)$, in which $a$ is
the least length for large $n$. Since $\Phi'(T_*)=0$, \cref{src:lem:gramcount} gives
$N_\Lambda(R)=1+2\Phi_*R^2+O(1)=n+O(1)$. On row one the boundary in the integer coordinate is at
$m_0=(R/a)\sqrt{1-T_*^{-2}}$, and moving $m$ a bounded number of steps past $m_0$ changes
$a^2m^2+a^{-2}$ by $O(1)$, since $aR=T_*$; so the discrepancy $O(1)$ can be corrected with each selected
or discarded point within $O(1)$ of squared radius $R^2$. The resulting set satisfies
$\sum_{v\in S}\abs v^2=nR^2-P_\Lambda(R)+O(1)=(n-1)^2/(4\Phi_*)+O(1)$. Hence
$\sqrt{E(n)}=(n-1)^2/(8\Phi_*)+O(1)$, and squaring gives \eqref{src:eq:gramasym}.

For the structure, isotropize a minimizing configuration by \eqref{src:eq:gramtrace}; its sum of squared
lengths is $2\sqrt{E(n)}$. The gap in \cref{src:lem:grammax} and \eqref{src:eq:gramcrude} force $T\in I$
for large $n$, and the upper bound with \eqref{src:eq:gramduality} and \eqref{src:eq:gramsharp} gives
$R^4(\Phi_*-\Phi(T))=O(1)$, so $T=T_*+O(R^{-2})$ by uniqueness and $\Phi''(T_*)<0$. The nonnegative slack
in \eqref{src:eq:gramduality}, the sum of $R^2-\abs v^2$ over omitted interior points and of $\abs v^2-R^2$
over selected exterior points, is $O(1)$. So every point of squared length below $R^2-O(1)$ is selected and
none above $R^2+O(1)$, with an absolute constant. The horizontal representative is selected, and the rows
$j\ge7$ are excluded since $7/T_*>1$. Each of the rows $1,\dots,6$ meets the circle away from a tangency,
and a bounded change of the squared radius changes its half-width in integer coordinates by $O(1)$,
uniformly in the shear $\theta$. The interval count on row $j$ gives
$N_j=(2w_jR^2/T)\sqrt{1-j^2/T^2}+O(1)$, and replacing $T$ by $T_*$ costs $O(1)$, which gives
\eqref{src:eq:gramrows}; every coefficient there is positive.
\end{proof}

\begin{theorem}[Algebraic degrees of the extremal constants; \stU]\label{src:thm:gramdegree}
The field $K=\Q(T_*^2)$ has degree $48$ and equals $\Q(\kappa)$. Each of $T_*$ and $\Phi_*$ has degree
$96$ over $\Q$.
\end{theorem}

\begin{proof}
Put $u=T_*^2$ and $b_j=30w_j$, so $(b_1,\dots,b_6)=(30,15,20,15,24,10)$. With two parameters $s,t$ let
\[
 B_j=\bigl(b_j(4j^2-x)+s\mathbf1_{j=1}+tb_j(x-j^2)\bigr)\sqrt{x-j^2},
\]
and let $\Omega(x,s,t)$ be the product of $B_1+\sum_{j=2}^6\varepsilon_jB_j$ over the $32$ choices of signs
$\varepsilon_j$. Changing the sign of any radical permutes the factors, up to a common sign change of all
$32$ factors for the first radical; treating the radicals as independent formal square roots, each occurs
with even exponent, so $\Omega\in\Z[x,s,t]$ has degree at most $48$ in $x$. Let $P(x)=\Omega(x,0,0)$ and
let $H_s,H_t$ be the derivatives in $s,t$ at $s=t=0$. On the positive branch at $x=u$ the all-plus factor
vanishes by \eqref{src:eq:gramroot}, so $P(u)=0$.

An exact certificate shows that $\bar P=P\bmod53$ has degree $48$ and is irreducible: $P,H_s,H_t$ are
evaluated at $x=0,\dots,48$ over $\F_{53}$ and interpolated, the product over signs being formed by
multiplying successively by the conjugate under each of the last five sign changes (in the basis indexed by
subsets $I\subseteq\set{1,\dots,6}$ multiplication sends $I,J$ to $I\triangle J$ with the factor
$\prod_{j\in I\cap J}(x-j^2)$, and the product rule carries both derivatives), and in $\F_{53}[X]$
\[
 X^{53^{48}}\equiv X\pmod{\bar P},\qquad\gcd(\bar P,X^{53^d}-X)=1\quad(d=16,24).
\]
As $2,3$ are the prime divisors of $48$, this excludes every proper factor degree, and Gauss's lemma gives
$[K:\Q]=48$.

The root $u$ of $P$ is simple; every radical is analytic near this real point, so exactly one sign factor
vanishes there, and differentiating the product and cancelling the nonvanishing factors gives
\[
 \frac{H_t(u)}{H_s(u)}=\frac{\sum_{j=1}^6b_j(u-j^2)^{3/2}}{\sqrt{u-1}},\qquad
 \kappa=\frac{2025u^4H_s(u)^2}{256(u-1)H_t(u)^2},
\]
with both derivatives nonzero at $u$; so $\kappa\in K$. The leading coefficient of $P$ is a $53$-adic unit,
so reduction maps $\Z_{(53)}[u]$, with the displayed denominator inverted, to $\F_{53}[X]/(\bar P)$, where
the denominator is nonzero; the certificate verifies by cross multiplication that neither the $16$th nor
the $24$th power of Frobenius fixes the image of $\kappa$, which therefore has degree $48$ over $\F_{53}$. A
primitive integer relation for $\kappa$ of smaller degree would reduce to a nonzero relation for this image,
so $\Q(\kappa)=K$.

Let $\ell$ be the leading coefficient of $P$. Since the degree is even, the norms of $u$ and $u-1$ are
$P(0)/\ell$ and $P(1)/\ell$, and the certificate gives residues $50\bmod53$ and $42\bmod59$, quadratic
nonresidues at primes where $\ell$ is a unit. So neither $u$ nor $u-1$ is a square in $K$. Now $T_*=\sqrt u$
and $\Phi_*=2H_t(u)\sqrt{u-1}/(45u^2H_s(u))$; since $\Q(\Phi_*)$ contains $\kappa$ it contains $K$, and both
quadratic extensions are proper.
\end{proof}

The coefficients of $\bar P$ and the degree checks are recorded in the same output file.

For $2\le r\le n-2$ let $E_r(n)$ be the least value of $\det(AA^{\srctr})/D^2$ over integer $r\times n$
matrices $A$ with every maximal minor nonzero, where $D$ is the gcd of their absolute values (the
quantity of \cref{src:thm:gram} with $z=0$ when $n\ge5$), so $E_2(n)=E(n)$.

\begin{corollary}[Contraction and duality; \stU]\label{src:cor:gramrank}
For $2\le r\le n-2$,
\[
 \frac{E_r(n)}{\binom nr}\ge\max\Bigl\{\frac{E(n-r+2)}{\binom{n-r+2}2},\ \frac{E(r+2)}{\binom{r+2}2}\Bigr\},
 \qquad E_{n-2}(n)=E(n).
\]
Moreover $E_r(n)=E_{n-r}(n)$ for $2\le r\le n-2$.
\end{corollary}

For fixed $r\ge3$ the contraction bound is of order $n^{r+2}$, below the true order $n^{3r-2}$ given by
\cref{spn:thm:gramorder} below; the duality $E_r(n)=E_{n-r}(n)$ transfers that order to fixed corank.

\begin{proof}
Fix $r-2$ column indices $I$. A unimodular row operation carries the rational span of those columns to the
first $r-2$ coordinate axes, by a basis of its saturated lattice. Every minor extending $I$ is then, up to
sign, a common nonzero integer times a two-row minor of the projected columns; dividing this local family by
its own gcd gives normalized rank-two minors, and the global gcd divides the local one, so the sum of the
squared globally normalized minors extending $I$ is at least $E(n-r+2)$. Each full minor occurs $\binom r2$
times in the sum over $I$; dividing by the number of incidences gives the first bound. The Hodge star of the
primitive exterior product of a basis of the saturated rational row space is, up to sign, the primitive
exterior product for its orthogonal complement, since it spans the right line over $\Q$ and signed
complementary coordinates preserve primitivity. So $E_r(n)=E_{n-r}(n)$, and the first bound applied to
$n-r$ gives the second.
\end{proof}

\begin{example}[Finite witnesses]\label{src:ex:gramwitness}
Rectangular lattices with rational parameters close to $\Phi_*$ and $T_*$, with exact comparison of squared
lengths, give the following upper bounds for $E(n)$; the columns contain $(1,0)^{\srctr}$ and a point on row
one, so their index is one.
\begin{center}\small
\begin{tabular}{rrr}
\toprule
$n$ & $\det(AA^{\srctr})$ & $\det(AA^{\srctr})/(n-1)^4$\\
\midrule
$100$ & $6\,546\,225$ & $0.068147535$\\
$1\,000$ & $67\,885\,760\,907$ & $0.068157984$\\
$10\,000$ & $681\,306\,789\,878\,949$ & $0.068157938$\\
\bottomrule
\end{tabular}
\end{center}
The row counts of the witnesses are recorded with their Gram matrices in a delivered
file.\footnote{\raggedright\artifact{gram-euler-and-zero-filters/gram/examples.json}}
\end{example}

\paragraph{The order of growth in every fixed rank.} Call $A\in\Z^{r\times n}$ \emph{generic} if every $r$
of its columns are linearly independent, and for $A$ of rank $r$ write $\mathcal E(A)=\det(AA^{\srctr})/D^2$
with $D$ the gcd of the nonzero maximal minors, so that $E_r(n)$ is the least $\mathcal E(A)$ over generic
$A$. A homogeneous hyperplane is the kernel of a nonzero real linear functional. The lower bound rests on a
covering lemma that does not assume the box to be roughly cubic, so that it survives an arbitrary change of
lattice metric.

\begin{lemma}[Anisotropic box covering; \stU]\label{spn:lem:gramcover}
Let $r\ge2$, let $L_1,\dots,L_r\ge1$ and $P=\prod_iL_i$. The integer points of $\prod_i[-L_i,L_i]$ are covered
by at most $K_rP^{1/(r-1)}$ homogeneous hyperplanes with nonzero integer normals, where $K_r=(2r+3)^r$. In
particular, an indexed family of integer points in this box with at most $h$ members on each homogeneous
hyperplane has at most $hK_rP^{1/(r-1)}$ members.
\end{lemma}

\begin{proof}
Order the coordinates so that $L_1\le\dots\le L_r$, put $R_j=(L_1\cdots L_j)^{1/(j-1)}$ for $2\le j\le r$, and
let $j$ minimize $R_j$, $R=R_j$. Then $1\le R\le R_r=P^{1/(r-1)}$ and $R\ge L_j$: this is clear for $j=2$, and
for $j>2$, $L_j>R_j$ would give $R_{j-1}^{j-2}=R_j^{j-1}/L_j<R_j^{j-2}$. So $R/L_i\ge1$ for $i\le j$. Put
$C=r+1$, $m_i=\floor{CR/L_i}$ and $\mathcal B=\prod_{i\le j}\set{0,\dots,m_i}$. For an integer point $z$ of the box
the values $a\cdot(z_1,\dots,z_j)$, $a\in\mathcal B$, lie in an interval of length at most $jCR$, while
$\abs{\mathcal B}>(CR)^j/(L_1\cdots L_j)=C^jR\ge jCR+1$, since $C^j-jC\ge C(C-j)\ge3$ and $R\ge1$. Two distinct
$a,a'\in\mathcal B$ therefore give a nonzero integer normal $u=a-a'$, padded by zeros, with $u\cdot z=0$ and
$\abs{u_i}\le m_i$. These normals do not depend on $z$, and there are at most
$\prod_{i\le j}(2m_i+1)\le(2C+1)^jR^j/(L_1\cdots L_j)=(2C+1)^jR\le K_rP^{1/(r-1)}$ of them.
\end{proof}

\begin{lemma}[A short lattice basis; \stU]\label{spn:lem:gramshort}
Every full lattice $\Lambda\subset\R^r$ has a basis $u_1,\dots,u_r$ with
$\prod_i\norm{u_i}\le M_r\operatorname{covol}(\Lambda)$, where $M_r=(2/\sqrt3)^{r(r-1)/2}$.
\end{lemma}

\begin{proof}
Induct on $r$, the case $r=1$ being an equality. A shortest nonzero $u\in\Lambda$ is primitive and extends to a
basis, so the orthogonal projection $\Lambda'$ of $\Lambda$ to $u^\perp$ is a lattice with
$\operatorname{covol}(\Lambda)=\norm u\operatorname{covol}(\Lambda')$. Each $w\in\Lambda'$ has a lift $v$ whose
component along $u$ has length at most $\norm u/2$; then $\norm w^2+\norm u^2/4\ge\norm v^2\ge\norm u^2$ for
$w\ne0$, so $\norm w\ge(\sqrt3/2)\norm u$ and $\norm v\le(2/\sqrt3)\norm w$. Lifts of a basis of $\Lambda'$ given by
induction, together with $u$, form a basis of $\Lambda$, and $M_r=(2/\sqrt3)^{r-1}M_{r-1}$.
\end{proof}

\begin{theorem}[Gram bound with hyperplane occupancy; \stU]\label{spn:thm:gramrobust}
Let $A\in\Z^{r\times n}$ have rank $r\ge2$, and let $h$ be the largest number of its indexed columns lying in a
homogeneous hyperplane. If $h\le n/2$, then
\begin{equation}\label{spn:eq:gramrobust}
 \mathcal E(A)\ge\frac{n^r}{M_r^2(2r)^r}\Bigl(\frac n{2hK_r}\Bigr)^{2r-2}.
\end{equation}
\end{theorem}

\begin{proof}
Let $B$ be a basis matrix of the lattice $\Lambda_A$ generated by the columns, and write $A=BC$. Then $C$ is
integral and its columns generate $\Z^r$, so $CJ=I_r$ for an integral $J$, and Cauchy--Binet for $CJ$ shows that
the maximal minors of $C$ have gcd one. Each maximal minor of $A$ is $\det B$ times the corresponding minor of
$C$, so $\abs{\det B}=D$ and $\det(CC^{\srctr})=\mathcal E(A)$; hyperplane occupancies are unchanged. Replace $A$
by $C$. Let $G=AA^{\srctr}$ and $w_i=G^{-1/2}v_i$, so that $\sum_iw_iw_i^{\srctr}=I_r$ and $\sum_i\norm{w_i}^2=r$.
With $R=\sqrt{2r/n}$, more than $n/2$ columns, forming a set $I$, satisfy $\norm{w_i}\le R$, and since
$\abs I>h$ they span $\R^r$. The $w_i$ lie in $G^{-1/2}\Z^r$, whose dual lattice $G^{1/2}\Z^r$ has covolume
$\sqrt{\det G}$; take a basis $u_1,\dots,u_r$ of it as in \cref{spn:lem:gramshort}. For $i\in I$ the vectors
$z_i=(\langle u_k,w_i\rangle)_k$ are integral and lie in the box of half-widths $L_k=R\norm{u_k}$; each $L_k\ge1$,
since otherwise the $k$th coordinates of all $z_i$ would vanish, contradicting spanning. Their product is
$P\le(2r/n)^{r/2}M_r\sqrt{\det G}$. The map $w\mapsto z$ is linear and invertible, so each homogeneous hyperplane
contains at most $h$ of the $z_i$, and \cref{spn:lem:gramcover} gives $n/2<\abs I\le hK_rP^{1/(r-1)}$. Squaring
gives \eqref{spn:eq:gramrobust}, as $\det G=\mathcal E(A)$.
\end{proof}

\begin{lemma}[Short projective lift; \stU]\label{spn:lem:gramlift}
Let $p$ be prime, $r\ge2$, and $\ell\subset\F_p^r$ a line. Some $v\in\Z^r$ with
$\norm v_\infty\le H=\floor{p^{(r-1)/r}}$ reduces to a nonzero element of $\ell$.
\end{lemma}

\begin{proof}
The $(H+1)^r>p^{r-1}$ points of $\set{0,\dots,H}^r$ cannot have distinct images in $\F_p^r/\ell$; the
difference of two with equal image reduces into $\ell$, and is nonzero modulo $p$ since $H<p$.
\end{proof}

\begin{proposition}[Full-spark short columns; \stU]\label{spn:prop:gramcurve}
For every prime $p\ge n\ge r\ge2$ there is a generic $A\in\Z^{r\times n}$ with entries of absolute value at most
$H=\floor{p^{(r-1)/r}}$, and $\mathcal E(A)\le\det(AA^{\srctr})\le n^rH^{2r}\le n^rp^{2r-2}$.
\end{proposition}

\begin{proof}
For distinct $t_1,\dots,t_n\in\F_p$ apply \cref{spn:lem:gramlift} to the lines spanned by
$q(t_i)=(1,t_i,\dots,t_i^{r-1})^{\srctr}$. The lifted columns reduce to nonzero multiples of the $q(t_i)$, so every
maximal minor is a nonzero multiple of a Vandermonde determinant modulo $p$, hence nonzero. The trace of
$AA^{\srctr}$ is at most $nrH^2$, and the arithmetic--geometric mean inequality gives
$\det(AA^{\srctr})\le(nH^2)^r$.
\end{proof}

Merely lifting an affine moment curve coordinate by coordinate would cost too much; the freedom to rescale each
column modulo $p$ is what makes the columns short.

\begin{theorem}[Order of growth; \stU]\label{spn:thm:gramorder}
For every fixed $r\ge2$ and every $n\ge2r$,
\[
 c_rn^{3r-2}\le E_r(n)\le2^{2r-2}n^{3r-2},\qquad
 c_r=\bigl[M_r^2(2r)^r\bigl(2(r-1)K_r\bigr)^{2r-2}\bigr]^{-1},
\]
so $E_r(n)/\binom nr=\Theta_r(n^{2r-2})$. The lower bound holds for every generic matrix.
\end{theorem}

\begin{proof}
A generic matrix has $h=r-1\le n/2$, and \cref{spn:thm:gramrobust} gives the lower bound. Bertrand's postulate
gives a prime $n\le p<2n$, and \cref{spn:prop:gramcurve} gives $E_r(n)\le n^r(2n)^{2r-2}$.
\end{proof}

For $r=2$ this is the order of \cref{src:thm:gramasym}; for $r=3$ the minimum is of order $n^7$ and the mean
squared normalized minor of order $n^4$. The constants are not optimized. By \cref{src:cor:gramrank}, at fixed
corank $c\ge2$ one has $E_{n-c}(n)=E_c(n)=\Theta_c(n^{3c-2})$ for $n\ge2c$.

\begin{corollary}[A bounded number of vanishing minors; \stU]\label{spn:cor:gramzeros}
Fix $r\ge2$ and $z\ge0$. The least $\mathcal E(A)$ over $A\in\Z^{r\times n}$ of rank $r$ with at most $z$
vanishing maximal minors is $\Theta_{r,z}(n^{3r-2})$. More generally \eqref{spn:eq:gramrobust} holds with $h$
replaced by $H_z=r-1+(r!z)^{1/r}$ whenever $H_z\le n/2$.
\end{corollary}

\begin{proof}
A homogeneous hyperplane containing $h\ge r$ columns accounts for $\binom hr\ge(h-r+1)^r/r!$ vanishing minors, so
$h\le H_z$, also for $h<r$; apply \cref{spn:thm:gramrobust}. Generic matrices belong to the class, and
\cref{spn:thm:gramorder} gives the upper bound.
\end{proof}

The corollary does not determine the optimum with exactly $z$ vanishing minors, nor a sharp constant.

\paragraph{Generic column numbers.} For integers $\Delta\ge1$ and $r\ge2$ let $g(\Delta,r)$ be the largest
number of columns of a generic $A\in\Z^{r\times n}$ whose largest absolute maximal minor is exactly $\Delta$,
and $g_\le(\Delta,r)$ the same with largest absolute maximal minor at most $\Delta$. Kriepke, Kyureghyan and
Schymura prove $g(\Delta,r)=O_r(\Delta^{2/r})$ for $r\ge3$, construct generic matrices with about
$\Delta^{1/(r-1)}$ columns along the sequence $\Delta=\ceil{(r-1)^{(r-1)/2}}(p-1)^{r-1}$, $p$ prime
\cite[Proposition~3.3]{KriepkeKyureghyanSchymura2023Size}, and ask in their Question~3.4 whether $g(\Delta,r)=\Theta(\Delta^{1/(r-1)})$
for fixed $r\ge3$ \cite{KriepkeKyureghyanSchymura2023Size}. The answer is affirmative.

\begin{theorem}[Generic column numbers; \stU]\label{spn:thm:gramcolumn}
For every fixed $r\ge2$, as $\Delta\to\infty$,
\[
 g(\Delta,r)=\Theta_r(\Delta^{1/(r-1)}),\qquad g_\le(\Delta,r)=\Theta_r(\Delta^{1/(r-1)}).
\]
For every sufficiently large $\Delta$ there is a generic matrix with at least $b_r\Delta^{1/(r-1)}$ columns,
$b_r=2^{-r-1}(4r)^{-r}$, in which exactly one maximal minor, that of the first $r$ columns, has absolute value
$\Delta$, and all others are nonzero of absolute value at most $\Delta/2$.
\end{theorem}

\begin{lemma}[Balanced prescribed determinant; \stU]\label{spn:lem:gramanchor}
For all $\Delta\ge1$ and $r\ge2$ some $B\in\Z^{r\times r}$ has $\det B=\Delta$ and columns of length at most
$2\sqrt r\,\Delta^{1/r}$.
\end{lemma}

\begin{proof}
Put $b=\floor{\Delta^{1/r}}+1\le2\Delta^{1/r}$ and write $\Delta=\sum_{i<r}d_ib^i$ with $0\le d_i<b$. The columns
$be_j-e_{j+1}$, $1\le j\le r-1$, and $(d_0,\dots,d_{r-1})^{\srctr}$ have entries of absolute value at most $b$, and
expansion along the last column gives $\det B=\Delta$.
\end{proof}

\begin{lemma}[A prime avoiding the anchor; \stU]\label{spn:lem:gramprime}
Fix $r\ge2$ and $a>0$. For all large $\Delta$ some prime $p\nmid\Delta$ satisfies $p\ge r$ and
$2^{-r-1}a\Delta^{1/(r-1)}\le p\le a\Delta^{1/(r-1)}$.
\end{lemma}

\begin{proof}
Let $X=a\Delta^{1/(r-1)}$ and $m=\floor{X/2^r}$, so $m\ge\max(2,r)$ and $m\ge X/2^{r+1}$ for large $\Delta$. Bertrand's
postulate in $(2^im,2^{i+1}m)$, $0\le i<r$, gives $r$ distinct primes in $(m,X)$. Their product exceeds
$m^r>\Delta$ for large $\Delta$, as $r/(r-1)>1$, so one of them does not divide $\Delta$.
\end{proof}

\begin{proof}[Proof of \cref{spn:thm:gramcolumn}]
Upper bound: if a generic $A$ with $n\ge2r$ columns has all maximal minors at most $\Delta$, then
\cref{spn:thm:gramorder} and Cauchy--Binet give $c_rn^{3r-2}\le\mathcal E(A)\le\binom nr\Delta^2\le n^r\Delta^2$, so
$n\le c_r^{-1/(2r-2)}\Delta^{1/(r-1)}$; smaller $n$ are absorbed in the constant. Bounded version: the matrices of
\cref{spn:prop:gramcurve} with $n=p$ have maximal minors at most $r^{r/2}p^{r-1}$ by Hadamard's inequality, and
Bertrand's postulate chooses $p$. Exact version: take $a=(4r)^{-r}$ and $p$ from \cref{spn:lem:gramprime}, and $B$
from \cref{spn:lem:gramanchor}. Since $p\nmid\Delta$, $\bar B$ is invertible modulo $p$, and
$T=\bar B[q(0)\ \cdots\ q(r-1)]^{-1}\in\GL_r(\F_p)$ sends $q(0),\dots,q(r-1)$ to the columns of $\bar B$. Keep the
columns of $B$, and for each $t\in\F_p\setminus\set{0,\dots,r-1}$ append a lift of the line $\F_pTq(t)$ from
\cref{spn:lem:gramlift}. The $p$ columns reduce to nonzero multiples of the $Tq(t)$, so every maximal minor is
nonzero modulo $p$. Each appended column has length at most
$\sqrt r\,p^{(r-1)/r}\le\sqrt r(4r)^{1-r}\Delta^{1/r}\le\Delta^{1/r}/(2(2\sqrt r)^{r-1})$, the ratio of the
constants being $(2\sqrt r)^{2-r}\le1$, and is shorter than the anchor bound $2\sqrt r\,\Delta^{1/r}$. Every minor
other than $\det B=\Delta$ uses an appended column, and Hadamard's inequality bounds it by
$(2\sqrt r\,\Delta^{1/r})^{r-1}\Delta^{1/r}/(2(2\sqrt r)^{r-1})=\Delta/2$.
\end{proof}

No monotonicity of $g(\Delta,r)$ in $\Delta$ is used: the exact value is realized for every large $\Delta$, not
along a subsequence. For $r=3$ the exponent drops from the published $2/3$ to $1/2$. The statement is for fixed
rank and generic matrices; it is not uniform in $r$, says nothing about matrices with vanishing maximal minors, and
does not bear on the polynomial-time solvability of general $\Delta$-modular integer programs. Exact integer
checks of the constructions, from $(r,n,p)=(2,12,13)$ with $\mathcal E(A)=1949$ to $(5,23,23)$ with
$\mathcal E(A)=34\,155\,205\,413\,840$, all with gcd one, of the exact anchors $\Delta=100\,003$ for $r=2$ through
$\Delta=10^{18}+37$ for $r=5$, each attained by a single minor, and of the covering lemma on $17\,367$ points of seven
boxes, are recorded in
\artifact{gram_column_number_2026-11-27/data/verification.json}.

\paragraph{Rank two with many vanishing minors.} Now let $A=[v_1\ \cdots\ v_n]\in\Z^{2\times n}$ have rank two,
with zero, repeated and opposite columns allowed, let $z(A)$ be the number of pairs $i<j$ with
$\det(v_i,v_j)=0$, and let
\[
 E_\le(n,z)=\min_{z(A)\le z}\mathcal E(A),\qquad E_=(n,z)=\min_{z(A)=z}\mathcal E(A),
\]
the latter when the class is nonempty. With $w_j$ and $\Phi$ as in \eqref{src:eq:gramPhi} put
\begin{equation}\label{spn:eq:gramPsi}
 \Psi(T)=\frac{16}{15T}\sum_{1\le j\le T}w_j\Bigl(1-\frac{j^2}{T^2}\Bigr)^{5/2}\qquad(T>0),\qquad
 \Psi_*=\sup_{T>0}\Psi(T).
\end{equation}

\begin{theorem}[Sharp law with many vanishing minors; \stU]\label{spn:thm:gramdegenerate}
Let $n\to\infty$ and let $z=z(n)\ge0$ be integers with $z/n\to\infty$ and $z/n^2\to0$. Then the exact class is
nonempty for all large $n$, and
\[
 E_\le(n,z)\sim E_=(n,z)\sim\frac{n^6}{729\Psi_*^2z^2}=\frac{4n^6}{729\Psi_*^2Q^2}\,(1+o(1)),\qquad Q=n+2z .
\]
\end{theorem}

So for $z\asymp n^{1+\alpha}$, $0<\alpha<1$, the minimum has order $n^{4-2\alpha}$, a different regime from
$z=0$, where \cref{src:thm:gramasym} uses the first truncated potential; here a squared potential gives the
constant. For a lattice $\Lambda$ of covolume one, with $\mathcal P(\Lambda)$ and the reduced basis
$(a,0),(a\theta,1/a)$ as before \cref{src:lem:gramcount}, $T=aR$, put
$P_s(\Lambda,R)=\sum_{v\in\mathcal P(\Lambda)}(R^2-\abs v^2)_+^s$ for $s=1,2$, so that $P_1=P_\Lambda$.

\begin{lemma}[The squared potential; \stU]\label{spn:lem:gramsquared}
\begin{enumerate}[label=\textup{(\roman*)}]
\item Uniformly in $\Lambda$, as $R\to\infty$, $P_2(\Lambda,R)=R^6\Psi(T)+O(R^5\log(2R))$.
\item For fixed $T>0$, uniformly in the shear, $P_1=R^4\Phi(T)+O_T(R^2)$, $P_2=R^6\Psi(T)+O_T(R^4)$, and
$\#\set{v\in\mathcal P(\Lambda):\abs v<R}=O_T(R^2)$; this remains true after removing a fixed finite set of
directions, given in the integer coordinates, of lengths $O_T(R)$.
\item $\Psi(T)\to1/\pi$ as $T\to\infty$ and $\Psi(32/3)>31839/100000>1/\pi$, so $\Psi_*$ is attained at some
$T_0>1$, and every maximizer satisfies $\Phi(T_0)=\frac32\Psi_*$.
\end{enumerate}
\end{lemma}

\begin{proof}
(i) On row $j\le T$ put $\eta_j=(R^2-j^2/a^2)^{1/2}$; the function $f(x)=(\eta_j^2-x^2)_+^2$ has integral
$\frac{16}{15}\eta_j^5$ and variation $2\eta_j^4\le2R^4$. On a shifted grid a compactly supported function of bounded
variation sums to its integral divided by the spacing up to its variation, so M\"obius expansion of the coprimality
condition, as in the proof of \cref{src:lem:gramcount}, gives the main term $R^6\Psi(T)$ with row errors
$O(R^4\tau(j))$, which sum to $O(R^5\log(2R))$ as $T\le(4/3)^{1/4}R$; the horizontal direction contributes at most
$R^4$. (ii) For fixed $T$ only boundedly many rows and residue classes occur, and the same estimate gives errors
$O_T(R^2)$ and $O_T(R^4)$; the count is $O_T(R/a+1)$; removing finitely many directions changes $P_s$ by
$O_T(R^{2s})$. (iii) From $\sum_{j\le x}w_j=6x/\pi^2+O(\log(2x))$ and $\int_0^1(1-u^2)^{5/2}\,du=5\pi/32$, partial
summation gives the limit. With $s_j=1024-9j^2$ and $r_j=\floor{10^{12}\sqrt{s_j}}$,
\[
 \Psi(32/3)\ge\frac{16}{5\cdot32^6\cdot10^{12}}\sum_{j=1}^{10}\frac{\varphi(j)}js_j^2r_j>\frac{31839}{100000},
\]
a comparison of rationals, and $1/\pi<71/223<31839/100000$. So the supremum is attained in a compact subinterval of
$(1,\infty)$. The summands vanish to order $5/2$ at integers, so $\Psi$ is differentiable, and termwise
$4\Phi(T)=6\Psi(T)+T\Psi'(T)$; at a maximizer $\Psi'(T_0)=0$.
\end{proof}

\begin{lemma}[Rounded parabolic construction; \stU]\label{spn:lem:gramparabolic}
Fix a maximizer $T_0$ and a finite set $F$ of primitive integer directions. Let $N\to\infty$ and $Q\equiv N\pmod2$
with $N/Q\to0$ and $Q/N^2\to0$, put $Z=(Q-N)/2$, and let $\eta\to0$ with $(Q/N^2+N/Q)/\eta\to0$. For large $N$ there
is a rank-two integral matrix with $N$ nonzero primitive columns, in directions outside $F$, whose number $Z'$ of
vanishing minors satisfies $0\le Z-Z'=O(\eta Q)$, and which in a lattice of covolume one has squared-length sum
at most $(1+o(1))4N^3/(27Q\Psi_*)$; the constant depends only on $T_0$ and $F$.
\end{lemma}

\begin{proof}
Take $R^2=\frac4{9\Psi_*}\frac{N^2}Q(1+10\eta)$, $a=T_0/R$ and the rectangular lattice with basis $(a,0),(0,1/a)$.
For primitive directions $v\notin F$ let $f_v=(R^2-\abs v^2)_+$, $P_s=\sum_vf_v^s$, $c=N/P_1$ and
$\tilde m_v=cf_v$, so $\sum_v\tilde m_v=N$. By \cref{spn:lem:gramsquared}(ii) and $\Phi(T_0)=\frac32\Psi_*$,
\[
 \sum_v\tilde m_v^2=N^2\frac{P_2}{P_1^2}=\frac Q{1+10\eta}\bigl(1+O(Q/N^2)\bigr),
\]
which lies between $Q-O(\eta Q)$ and $Q-5\eta Q$ for large $N$. Round each $\tilde m_v$ down and then add one to
enough entries with positive fractional part to restore the sum $N$; there are $d=O(R^2)=O(N^2/Q)=o(N)$ entries,
each changed by less than one, so $\sum_vm_v^2$ moves by at most $2N+d=o(\eta Q)$. Thus $\sum m_v^2\le Q$ with deficit
$O(\eta Q)$, and $Z'=(\sum m_v^2-N)/2$. Before rounding the squared-length sum is
$c(R^2P_1-P_2)=N(R^2-P_2/P_1)=\frac13NR^2(1+o(1))=(1+o(1))4N^3/(27Q\Psi_*)$, and rounding changes it by
$O(dR^2)=o(NR^2)$. Two fixed allowed directions on row one have multiplicities of order $cR^2\asymp Q/N\to\infty$, so
the rank is two; coordinates in the lattice basis, of determinant one, give an integral matrix.
\end{proof}

\begin{proof}[Proof of \cref{spn:thm:gramdegenerate}]
\emph{Lower bound.} As in the normalization before \eqref{src:eq:gramtrace}, a minimizer may be taken with columns
generating $\Z^2$ and every nonzero column primitive: passing to a basis of the column lattice and dividing a
column by its content keep the zero pattern and do not increase the determinant. If there are $k$ zero columns and
the nonzero directions have multiplicities $m_v$, then $\sum_vm_v=n-k$ and
$\sum_vm_v^2=n+2z(A)-2kn+k^2\le Q$, since the zero columns account for $k(n-k)+\binom k2\ge k(n-1)/2$ vanishing
minors; in particular $k\le2z/(n-1)=o(n)$. Isotropize by \eqref{src:eq:gramtrace}; the squared-length sum is
$S=2\sqrt{\mathcal E(A)}$, and for every $R>0$, by Cauchy--Schwarz,
\[
 S\ge(n-k)R^2-\sum_vm_v(R^2-\abs v^2)_+\ge(n-k)R^2-\sqrt{Q\,P_2(\Lambda,R)} .
\]
With $R=2(n-k)/(3\sqrt{Q\Psi_*})\to\infty$ and \cref{spn:lem:gramsquared}(i), $P_2\le R^6(\Psi_*+o(1))$ uniformly in
$\Lambda$, so $S\ge(1-o(1))4n^3/(27Q\Psi_*)$ and $E_\le(n,z)\ge(1-o(1))4n^6/(729Q^2\Psi_*^2)$. This uses only
$z=o(n^2)$, and bounds the exact class as well.

\emph{Upper bound with exactly $z$ vanishing minors.} Put $e=z/n^2+n/z\to0$, $\eps=e^{1/4}$, $M=\ceil{\sqrt{\eps z}}$, and
reserve the eight directions $(u,1)$, $0\le u\le7$, as $F$. Let $N=n-8M$, $Z=z-4M(M-1)$ and $Q'=N+2Z$; then $M=o(n)$,
$N\sim n$ and $Q'\sim2z\sim Q$. \Cref{spn:lem:gramparabolic} with $\eta=e^{1/2}$ applies, since
$Q'/N^2+N/Q'=O(e)$, and leaves an unused count $D=Z-Z'=O(e^{1/2}z)=o(M^2)$, so $D\le M^2$ for large $n$. Write
$D=x_1^2+\dots+x_4^2$ with $\abs{x_i}\le M$ by Lagrange's four-square theorem, and give the reserved directions
$2i-2,2i-1$ the multiplicities $M+x_i$ and $M-x_i$. They add $8M$ columns and
$\sum_i\bigl[\binom{M+x_i}2+\binom{M-x_i}2\bigr]=4M(M-1)+D$ vanishing minors, and distinct directions give none, so
the matrix has $n$ columns and exactly $z$ vanishing minors. The reserved vectors have squared length $O(R^2)$ with
$R^2\asymp n^2/z$, so they add $O(MR^2)=o(n^3/Q)$ to the squared-length sum. Since $\det M'\le(\tr M'/2)^2$ for a
positive semidefinite $2\times2$ matrix $M'$ and a gcd above one only lowers $\mathcal E$,
$E_=(n,z)\le(1+o(1))4n^6/(729Q^2\Psi_*^2)$. Finally $E_\le\le E_=$.
\end{proof}

In the construction the multiplicity of a direction $(a(u+j\theta),j/a)$ is approximately
$\frac{3Q}{2n}(1-x^2-j^2/T_0^2)_+$ with $x=T_0(u+j\theta)/R^2$, parabolic up to the boundary of an ellipse; about
$n^2/z$ directions are used, with typical multiplicity of order $z/n$, and only the rows $j<T_0$. Numerically the
maximum appears to be at $T\approx10.67368$ with $\Psi_*\approx0.3183946273$, giving
$1/(729\Psi_*^2)\approx0.01353135$, but neither uniqueness of the maximizer nor this value is certified; exact
evaluation gives $0.318394611593156<\Psi(32/3)<0.318394611593166$. Rounded parabolic profiles at $T=32/3$ give
integral witnesses with $\mathcal E(A)z^2/n^6$ equal to $0.0123790$ at $(n,z)=(4841,158\,600)$ and $0.0135188$ at
$(n,z)=(47\,765\,579,\,159\,199\,178\,481)$, and the four-square correction with $1357=0^2+3^2+18^2+32^2$, $M=40$,
gives a matrix with exactly $18\,048$ columns and $1\,140\,772$ vanishing minors:
\begin{center}\small
\artifact{proveit_three_extensions_2026-11-27_1/papers/gram_determinants/}\\
\artifact{data/verification.json}
\end{center}

\subsection{Every sufficiently large bound in a self-dual projective lattice}\label{spn:sec:projective}
\sources{grh_nonabelian_transfer_2026-11-25}

Let $K$ be imaginary quadratic, $\OO=\OO_K$, and assume that some
$\eta\in\OO$ satisfies
\begin{equation}\label{spn:eq:traceone}
 \eta+\overline\eta=1 .
\end{equation}
For the source, $K=\Q(\sqrt{-d})$ with $d\equiv7\pmod8$, and
$\eta=(1+\sqrt{-d})/2$ works. Let $V$ have positive dimension over $K$, and
let $h$ be a positive definite Hermitian form, conjugate-linear in the first
variable. For $v\in V$, define the rank-one operator $\Theta_v(z)=v\,h(v,z)$;
for $x,y\in V$, also put $\Theta_{x,y}(z)=x\,h(y,z)$. For a full
$\OO$-lattice $A\subset V$ write
$A^\#=\set{z\in V:h(z,A)\subseteq\OO}$.

\begin{theorem}[Projective Hermitian eventual frames; \stU]\label{spn:thm:projective}
Let $A$ be a full projective $\OO$-lattice in $V$ with $A^\#=A$. Let
$S\subset A$ be finite and $R\in\Z_{\ge0}$. There are a fixed finite family
$T\subset A$ containing $S$ and an integer $M_0$ such that for every integer
$M\ge M_0$ there are integers $m_v\ge R$ satisfying
\begin{equation}\label{spn:eq:projective-frame}
 \sum_{v\in T}m_v\Theta_v=M\,I_V .
\end{equation}
In orthonormal complex coordinates this is an ordinary tight-frame identity.
Neither freeness of $A$ nor a principal fractional ideal is assumed.
\end{theorem}

\begin{lemma}[An integral signed resolution of the identity]\label{spn:lem:signed}
Under the hypotheses of \cref{spn:thm:projective}, the identity is an integer
linear combination of finitely many $\Theta_v$ with $v\in A$.
\end{lemma}

\begin{proof}
The dual-basis lemma for a finite projective module supplies $x_i\in A$ and
$\OO$-linear functionals $\ell_i\colon A\to\OO$ such that
$z=\sum_ix_i\ell_i(z)$. Each functional extends to $V$ and, by
nondegeneracy, has the form $\ell_i(z)=h(y_i,z)$ for some $y_i\in V$. Since
$\ell_i(A)\subseteq\OO$, self-duality implies $y_i\in A$. Thus
\[
 \sum_i\Theta_{x_i,y_i}=I,\qquad\sum_i\Theta_{y_i,x_i}=I,
\]
where the second equality is the adjoint of the first. Direct expansion gives
the exact polarization identity
\begin{equation}\label{spn:eq:polarization}
 \Theta_{x+\eta y}-\Theta_x-\Nm_{K/\Q}(\eta)\,\Theta_y
 =\overline\eta\,\Theta_{x,y}+\eta\,\Theta_{y,x}.
\end{equation}
Summing over $i$ and applying \eqref{spn:eq:traceone} gives $I$ on the right.
The left is an integer linear combination of rank-one operators of lattice
vectors, since $\Nm_{K/\Q}(\eta)\in\Z$. Combine repeated vectors to obtain
\begin{equation}\label{spn:eq:signed-id}
 I=\sum_{v\in S'}z_v\Theta_v,\qquad z_v\in\Z .
\end{equation}
The prescribed set $S$ can be added to $S'$ with zero coefficients.
\end{proof}

The trace-one element removes the factor two that a purely real polarization
argument would introduce. Projective dual generators replace a global integral
basis; this is why no freeness assumption enters the proof.

\begin{lemma}[A positive reserve containing prescribed vectors]\label{spn:lem:reserve}
For every finite $S'\subset A$, there are a positive integer $q$, a finite
$T\supseteq S'$, and positive integers $c_v$ such that
\begin{equation}\label{spn:eq:reserve}
 q\,I=\sum_{v\in T}c_v\Theta_v .
\end{equation}
\end{lemma}

\begin{proof}
Choose a sufficiently small positive rational $\epsilon$ so that
$I-\epsilon\sum_{v\in S'}\Theta_v$ is positive definite. Choose a $K$-basis of
$V$ inside $A$, with complex coordinate matrix $B$. Its Hermitian Gram matrix
$H=B^*B$ has entries in $K$. If $u_v\in K^n$ is the coordinate column of $v$,
then
\[
 I=BH^{-1}B^*,\qquad\Theta_v=Bu_vu_v^*B^* .
\]
The residual matrix $H^{-1}-\epsilon\sum_{v\in S'}u_vu_v^*$ is positive
definite Hermitian over $K$. Hermitian $LDL^*$ decomposition has a triangular
factor over $K$ and strictly positive diagonal entries in $K\cap\R=\Q$. This
expresses the residual as a sum of positive rational multiples of rank-one
matrices with vectors in $K^n$.

Multiply each vector by an integer so that its image in $V$ lies in $A$, and
divide its coefficient by the corresponding square. Reinsert the vectors of
$S'$ with coefficient $\epsilon$. Clearing all positive rational denominators
gives \eqref{spn:eq:reserve}. Zero vectors, if prescribed, may be retained
with arbitrary positive coefficients. Combine duplicate vectors if necessary.
\end{proof}

\begin{proof}[Proof of \cref{spn:thm:projective}]
Use \eqref{spn:eq:signed-id} and \eqref{spn:eq:reserve}, extending $z_v=0$
outside the signed support. Define
\begin{equation}\label{spn:eq:explicit-threshold}
 k_0=\max_{v\in T}\ceil*{\frac{R+(q-1)\max(0,-z_v)}{c_v}},\qquad M_0=qk_0 .
\end{equation}
For $M\ge M_0$, write $M=kq+r$ with $0\le r<q$, and set $m_v=kc_v+rz_v$. Then
$k\ge k_0$ and every $m_v\ge R$. Moreover,
\[
 \sum_{v\in T}m_v\Theta_v=k\sum_{v\in T}c_v\Theta_v+r\sum_{v\in T}z_v\Theta_v
 =kqI+rI=MI .
\]
This proves the theorem with a fixed family independent of $M$.
\end{proof}

\begin{corollary}[Coprime and prescribed large bounds; \stU]\label{spn:cor:coprime}
Let $A$ be as in \cref{spn:thm:projective}. Then $A$ has frames of every
sufficiently large positive integer bound; in particular it has coprime pairs
of bounds, and arbitrarily large bounds prime to any fixed integer, and any
finite set of its elements can be prescribed as frame members with any fixed
lower multiplicity. This applies to the corrected ideal $A=A_0\OO_E$ of the
unramified nonabelian benchmark of \cref{grp:sec:s3}, with both unit witnesses
inherited from the source prescribed as frame members.
\end{corollary}

\begin{proof}
Consecutive sufficiently large integers give coprime bounds; integers
congruent to one modulo a fixed integer give the other assertion. The
corrected ideal is Hermitian self-dual and projective, and $\OO_K$ contains
the required trace-one element.
\end{proof}

The inherited two-element witness set is unusually convenient. A qualitative
substitute exists for any fractional ideal.

\begin{lemma}[Finite simultaneous local generators; \stU]\label{spn:lem:local-generators}
For a nonzero fractional ideal $A$ of a number field $E$, there is a finite
set $S\subset A$ such that, for every rational prime $p$, some $y\in S$
satisfies
\[
 v_{\mathfrak P}(y)=v_{\mathfrak P}(A)\quad\text{for all }\mathfrak P\mid p .
\]
\end{lemma}

\begin{proof}
Choose $0\ne y_0\in A$. The integral ideal $(y_0)A^{-1}$ has finite support.
Outside the finitely many rational primes below that support, $y_0$ already
works. For each exceptional $p$, choose a nonzero class in every
one-dimensional residue module $A/\mathfrak PA$, for $\mathfrak P\mid p$. The
Chinese remainder theorem gives one element $y_p\in A$ representing those
classes simultaneously. Its local valuation is exactly that of $A$ at every
such prime. Take $S$ to consist of $y_0$ and these finitely many $y_p$.
\end{proof}

\begin{remark}[Limits of the qualitative theorem]\label{spn:rem:projective-limits}
For each fixed tower level, qualitative coprime frames and prescribed
sufficiently large bounds are available in the correct lattice; the raw
integral obstruction of \cref{spn:thm:dual-obstruction} is overcome by the
fractional-lattice construction, not contradicted. The threshold
\eqref{spn:eq:explicit-threshold} may be very large. Rational decomposition
and denominator clearing do not by themselves bound $q$, $c_v$, $z_v$, the
coordinates of the vectors, or $M_0$ uniformly in a growing tower. The source
needs a small auxiliary bound $h\le b$ and an exact final bound $M=p_0p_1$ in
a prescribed window, with multiplicities at least $6fb$ and controlled
conjugate norms. None of these quantitative inequalities follows from
qualitative eventual existence.
\end{remark}

%======================================================================
\subsection{Open questions}\label{spn:sec:open}

\begin{question}[Optimal selection]\label{spn:q:selection}
Under \eqref{spn:eq:near} and \eqref{spn:eq:moment}, the possible total ranks for a fixed quadruple
$(d,s,u,v)$ are determined by \cref{spn:thm:rank-spectrum}, but the optimal number of translates
needed to reach a given intermediate rank is not.
\end{question}

\begin{question}[Cheap completions]\label{spn:q:completion}
For a target beyond the source, construct integral completing vectors with
excess completion cost $K-A=O(E_0)$ in \eqref{spn:eq:rank-error}, with control
of their twist scalars and heights, so that a fixed positive fraction of
independent short forms suffices.
\end{question}

\begin{question}[Quantitative frames]\label{spn:q:frames}
Bound the reserve denominator $q$, the signed coefficients $z_v$, the vector
sizes and the threshold $M_0$ of \cref{spn:thm:projective} uniformly as the
tower modulus grows; a polynomial bound in the modulus with fixed exponents
would be a very favorable outcome. The source needs a small auxiliary bound
$h\le b$, an exact final bound $M=p_0p_1$ in a prescribed window,
multiplicities at least $6fb$ and controlled conjugate norms, and a threshold
growing like $\exp(f^{c})$ for a fixed $c>0$ would exceed the permitted
scale.
\end{question}

\begin{question}\label{src:q:gram}
For fixed $(r,n,z)$, what is the least value of $\det(AA^{\srctr})/D^2$ in
\cref{src:thm:gram}? The three affine bounds are each attained for $r=2$, $n=5$, but their
maximum is not claimed to be optimal for every triple. For $r=2$ and $z=0$,
\cref{src:thm:gramasym} gives $E(n)=\kappa(n-1)^4+O(n^2)$: what are the limit points of
$(E(n)-\kappa(n-1)^4)/n^2$, at which the row endpoints and the basis shear enter, and from which $n$
on are all minimizers confined to six rows? For $r=2$ and a fixed number of vanishing minors, where
the multiplicities of repeated directions must also be optimized, \cref{spn:cor:gramzeros} gives the
order $n^4$ when at most $z$ minors vanish: what is the asymptotic minimum, with at most or with exactly
$z$ vanishing minors? For $n\ll z\ll n^2$, \cref{spn:thm:gramdegenerate} gives the constant
$1/(729\Psi_*^2)$: is the maximizer of $\Psi$ unique, which rows do minimizers occupy, and what happens
when $z$ is of order $n$ or of order $n^2$? For fixed $r\ge3$, \cref{spn:thm:gramorder} gives
$E_r(n)=\Theta_r(n^{3r-2})$: does $E_r(n)/n^{3r-2}$ converge, to which constant, and what do minimizers
look like? What are the exact values of the generic column numbers $g(\Delta,r)$ of
\cref{spn:thm:gramcolumn}, and their sharp constants?
\end{question}
